% =====================================================================
%  Báo cáo Tiểu luận 01 — Phát triển các Hệ thống Thông minh
%  Nguyễn Duy Nghĩa · D23CTPM01 · B23DCCN600
%
%  Dựng bằng một lệnh:  python build.py
%  build.py sinh macro số liệu, bảng, trích mã nguồn rồi gọi
%      latexmk -lualatex -shell-escape main.tex
%  Hình dữ liệu là PDF vector vẽ bằng matplotlib (figures/); sơ đồ lý thuyết là TikZ (tikz/),
%  được externalize vào tikzcache/.
% =====================================================================
% extreport chỉ có 8-12, 14, 17, 20pt: khai báo 13pt bị bỏ qua và cả bài rơi về 10pt. Dùng 14pt.
\documentclass[a4paper,14pt,oneside]{extreport}

% --------------------------- Ngôn ngữ & font -------------------------
\usepackage{fontspec}
\setmainfont{Times New Roman}
\setsansfont{Segoe UI}
\directlua{luaotfload.add_fallback("toanhoc", {"Cambria Math:mode=node;"})}
\setmonofont{Consolas}[Scale=0.9, RawFeature={fallback=toanhoc}]
\usepackage{polyglossia}
\setmainlanguage{vietnamese}
\gappto\captionsvietnamese{%
  \renewcommand{\listfigurename}{Danh mục Hình vẽ}%
  \renewcommand{\listtablename}{Danh mục Bảng biểu}%
  \renewcommand{\contentsname}{Mục lục}}

% --------------------------- Khổ trang -------------------------------
\usepackage[a4paper,left=3cm,right=2cm,top=2.5cm,bottom=2.5cm,headheight=20pt]{geometry}
\usepackage{setspace}
\onehalfspacing
\setlength{\emergencystretch}{3em}   % cho phép giãn khoảng trắng thêm thay vì để từ dài (tên lớp, \texttt) tràn lề

% --------------------------- Toán học --------------------------------
\usepackage{amsmath,amssymb,amsfonts}
\usepackage{bm}

% --------------------------- Hình & bảng -----------------------------
\usepackage{graphicx}
\graphicspath{{assets/}{figures/}}
\usepackage{booktabs,longtable,array,tabularx,multirow}
\usepackage{float}
\usepackage[table,dvipsnames]{xcolor}
\usepackage[font=normalsize,labelfont=bf,labelsep=period]{caption}
\usepackage{subcaption}
\usepackage{adjustbox}
\usepackage{makecell}   % tiêu đề cột hai dòng trong bảng do build.py sinh
\renewcommand{\cellset}{\renewcommand{\arraystretch}{0.95}}

% --------------------------- Danh sách -------------------------------
\usepackage{enumitem}
\setlist[itemize]{topsep=3pt,itemsep=2pt,parsep=0pt,leftmargin=1.5em}
\setlist[enumerate]{topsep=3pt,itemsep=2pt,parsep=0pt,leftmargin=1.7em}

% --------------------------- Màu & tiêu đề ---------------------------
\definecolor{ptitblue}{HTML}{123A75}
\definecolor{ptitred}{HTML}{C00000}
\definecolor{softgray}{HTML}{F2F4F7}

\usepackage{titlesec}
\titleformat{\chapter}[display]
  {\normalfont\sffamily\bfseries\color{ptitblue}}
  {\Large\chaptertitlename\ \thechapter}{6pt}{\LARGE}
\titlespacing*{\chapter}{0pt}{-18pt}{22pt}
\titleformat{\section}{\normalfont\large\sffamily\bfseries\color{ptitblue}}{\thesection}{0.6em}{}
\titleformat{\subsection}{\normalfont\sffamily\bfseries\color{ptitblue}}{\thesubsection}{0.6em}{}
\titleformat{\subsubsection}{\normalfont\bfseries\itshape}{\thesubsubsection}{0.6em}{}
\setcounter{secnumdepth}{3}
\setcounter{tocdepth}{2}

% --------------------------- Đầu trang / chân trang ------------------
\usepackage{fancyhdr}
\pagestyle{fancy}
\fancyhf{}
\renewcommand{\headrulewidth}{0.4pt}
\fancyhead[L]{\small\itshape Nhóm lớp: CT \textbar{} Nhóm: 07 -- Nguyễn Duy Nghĩa}
\fancyhead[R]{\small\itshape Tiểu luận: Phát triển Hệ thống Thông minh}
\fancyfoot[C]{\thepage}
\fancypagestyle{plain}{\fancyhf{}\fancyfoot[C]{\thepage}\renewcommand{\headrulewidth}{0pt}}

% --------------------------- Tài liệu tham khảo ----------------------
\usepackage[backend=biber,style=numeric,sorting=none,maxbibnames=6]{biblatex}
\DeclareLanguageMapping{vietnamese}{english}
\addbibresource{refs.bib}

% --------------------------- TikZ / pgfplots -------------------------
% =====================================================================
%  Style dùng chung cho mọi hình TikZ/pgfplots của báo cáo.
%  Bảng màu phong phú, khối có bóng đổ và chuyển sắc, macro khối hộp 3D.
% =====================================================================
\usepackage{tikz}
\usepackage{pgfplots}
\pgfplotsset{compat=1.18}
\usetikzlibrary{external,positioning,arrows.meta,matrix,fit,calc,shapes.geometric,shapes.misc,shapes.arrows,
                backgrounds,decorations.pathreplacing,decorations.pathmorphing,shadows,shadows.blur,patterns,3d,
                perspective,fadings,shadings,trees}
\usepgfplotslibrary{groupplots,colormaps,fillbetween}

% Externalize: mỗi hình biên dịch một lần thành tikzcache/<tên>.pdf.
\tikzexternalize[prefix=tikzcache/]
\newcommand{\hinh}[1]{\tikzsetnextfilename{#1}\input{tikz/#1}}

% ----------------------------------------------------------------- bảng màu
\definecolor{mZero}{HTML}{7F7F7F}
\definecolor{mOne}{HTML}{0072B2}
\definecolor{mTwo}{HTML}{E69F00}
\definecolor{mThree}{HTML}{009E73}
\definecolor{accentRed}{HTML}{C0392B}
\definecolor{lightBlue}{HTML}{DCE8F5}
\definecolor{lightOrange}{HTML}{FCE9CF}
\definecolor{lightGreen}{HTML}{D5F0E6}
\definecolor{lightGray}{HTML}{EDEDED}
% màu tươi cho sơ đồ
\definecolor{cBlue}{HTML}{2F80ED}
\definecolor{cSky}{HTML}{56CCF2}
\definecolor{cTeal}{HTML}{1ABC9C}
\definecolor{cGreen}{HTML}{27AE60}
\definecolor{cLime}{HTML}{A3D977}
\definecolor{cYellow}{HTML}{F2C94C}
\definecolor{cOrange}{HTML}{F2994A}
\definecolor{cRed}{HTML}{EB5757}
\definecolor{cPink}{HTML}{E86FA8}
\definecolor{cPurple}{HTML}{9B51E0}
\definecolor{cIndigo}{HTML}{5B5FE0}
\definecolor{cBrown}{HTML}{A0522D}
\definecolor{cNavy}{HTML}{123A75}
\definecolor{cSlate}{HTML}{4F5D75}

\pgfplotsset{
  a06/.style={
    width=8cm, height=5.6cm, grid=major, grid style={gray!25},
    tick label style={font=\small}, label style={font=\normalsize},
    legend style={font=\small, draw=gray!50, fill=white, fill opacity=0.9, text opacity=1, rounded corners=2pt},
    title style={font=\normalsize\bfseries},
    /pgf/number format/use comma, /pgf/number format/1000 sep={.}},
  m0/.style={mZero, very thick, mark=none},
  m1/.style={mOne, very thick, mark=none},
  m2/.style={mTwo, very thick, mark=none},
  m3/.style={mThree, very thick, mark=none},
}

\tikzset{
  % khối chữ nhật bo góc có bóng mờ và chuyển sắc từ trên xuống
  hop/.style={rounded corners=3pt, draw=#1!75!black, thick, top color=#1!12, bottom color=#1!38,
              blur shadow={shadow blur steps=6, shadow xshift=0.8pt, shadow yshift=-0.8pt, shadow opacity=25},
              align=center, inner sep=4pt, font=\small},
  hop/.default=cBlue,
  khoi/.style={hop=cBlue, minimum height=8mm},
  khoiC/.style={hop=cOrange, minimum height=8mm},
  khoiG/.style={hop=cGreen, minimum height=8mm},
  khoiX/.style={hop=cSlate, minimum height=8mm},
  khoiP/.style={hop=cPurple, minimum height=8mm},
  khoiR/.style={hop=cRed, minimum height=8mm},
  khoiT/.style={hop=cTeal, minimum height=8mm},
  % quả cầu có bóng (nơ-ron, phép toán)
  cau/.style={circle, inner color=white, outer color=#1!55, draw=#1!65!black, thick, text=black!85,
              font=\small\bfseries, inner sep=1pt, minimum size=8mm,
              blur shadow={shadow blur steps=5, shadow opacity=30}},
  cau/.default=cBlue,
  mui/.style={-{Stealth[length=2.6mm]}, very thick, cNavy},
  muiM/.style={-{Stealth[length=2.6mm]}, very thick, #1},
  nhan/.style={font=\small, text=black!75},
  o/.style={draw=gray!70, minimum size=7mm, inner sep=0pt, font=\small},
  tieude/.style={font=\normalsize\bfseries\sffamily, text=cNavy},
}

% \khoibaD{x}{y}{rộng}{cao}{sâu}{màu}: khối hộp 3D, góc dưới-trái-trước tại (x,y),
% chiều sâu vẽ xiên lên phải (0,6 × sâu theo 30 độ). Mặt trên sáng, mặt bên tối.
\newcommand{\khoibaD}[6]{%
  \pgfmathsetmacro{\kdx}{0.6*#5*cos(30)}\pgfmathsetmacro{\kdy}{0.6*#5*sin(30)}%
  \filldraw[draw=#6!60!black, fill=#6!85!black, line join=round] (#1,#2) rectangle ++(#3,#4);
  \filldraw[draw=#6!60!black, fill=#6!45, line join=round] (#1,#2+#4) -- ++(\kdx,\kdy) -- ++(#3,0) -- ++(-\kdx,-\kdy) -- cycle;
  \filldraw[draw=#6!60!black, fill=#6!60!black, line join=round] (#1+#3,#2) -- ++(\kdx,\kdy) -- ++(0,#4) -- ++(-\kdx,-\kdy) -- cycle;
}


% --------------------------- Mã nguồn --------------------------------
\usepackage[most]{tcolorbox}
\tcbuselibrary{external}
\tcbset{shield externalize}
\usepackage{minted}
\setminted{fontsize=\small,linenos,numbersep=6pt,breaklines,frame=lines,
           framesep=2mm,baselinestretch=1.05,tabsize=4}
\usemintedstyle{friendly}
\renewcommand{\listingscaption}{Đoạn mã}
\renewcommand{\listoflistingscaption}{Danh mục Đoạn mã nguồn}
\counterwithin{listing}{chapter}
% \maNguon{module__ten}{chú thích}: đoạn mã build.py trích nguyên văn từ src/TieuLuan 01/tl/<module>.py
% (hoặc app/app.py), hàm hoặc lớp <ten>; mã trong báo cáo vì vậy luôn là mã đã chạy ra số liệu.
\newcommand{\maNguon}[2]{%
  \begin{listing}[H]
  \inputminted{python}{code/#1.py}
  \caption{#2}
  \label{code:#1}
  \end{listing}}

% --------------------------- Liên kết --------------------------------
\usepackage{hyperref}
\hypersetup{
  colorlinks=true, linkcolor=ptitblue, citecolor=ptitblue,
  urlcolor=ptitblue, filecolor=ptitblue,
  pdftitle={Tiểu luận: Học máy cơ bản, CNN và RNN trong hệ thống thông minh},
  pdfauthor={Nguyễn Duy Nghĩa — B23DCCN600 — Nhóm lớp CT — Nhóm 07},
  pdfsubject={Phát triển các Hệ thống Thông minh},
  bookmarksnumbered=true
}
\usepackage{bookmark}
\usepackage{microtype}

% --------------------------- Lệnh riêng ------------------------------
\usepackage{etoolbox}
\newcommand{\nhanxet}[1]{\par\smallskip\noindent\textbf{Nhận xét:} #1\par\smallskip}
\newtcolorbox{phathien}[1]{
  breakable,
  colback=softgray, colframe=ptitblue, boxrule=0.7pt, arc=2pt,
  left=8pt, right=8pt, top=6pt, bottom=6pt,
  fonttitle=\bfseries\sffamily, title={#1}
}
% \hinhpdf{tệp}{rộng}{chú thích}{nhãn}: hình dữ liệu matplotlib trong figures/.
\newcommand{\hinhpdf}[4]{%
  \begin{figure}[H]\centering
  \IfFileExists{figures/#1.pdf}{\includegraphics[width=#2\linewidth]{#1}}%
    {\fbox{\parbox[c][3cm][c]{0.8\linewidth}{\centering chưa vẽ: \detokenize{#1}}}\PackageWarning{solieu}{Thieu so lieu 'fig/#1'}}
  \caption{#3}\label{fig:#4}
  \end{figure}}
% \haihinh{tệp1}{chú thích 1}{tệp2}{chú thích 2}{chú thích chung}{nhãn}: hai hình nửa trang cạnh nhau (a), (b).
\newcommand{\haihinh}[6]{%
  \begin{figure}[H]\centering
  \begin{subfigure}[t]{0.49\linewidth}\centering\includegraphics[width=\linewidth]{#1}\caption{#2}\end{subfigure}\hfill
  \begin{subfigure}[t]{0.49\linewidth}\centering\includegraphics[width=\linewidth]{#3}\caption{#4}\end{subfigure}
  \caption{#5}\label{fig:#6}
  \end{figure}}
% \hinhtikz{tệp}{chú thích}{nhãn}: sơ đồ TikZ trong tikz/, co theo bề rộng trang nếu cần.
\newcommand{\hinhtikz}[3]{%
  \begin{figure}[H]\centering
  \begin{adjustbox}{max width=\linewidth}\hinh{#1}\end{adjustbox}
  \caption{#2}\label{fig:#3}
  \end{figure}}
% \anhUI{tệp}{chú thích}{nhãn}: ảnh chụp Web App (assets/ui/<tệp>.png) do app/screenshots.py sinh.
\newcommand{\anhUI}[3]{%
  \begin{figure}[H]\centering
  \IfFileExists{assets/ui/#1.png}{\fbox{\includegraphics[width=0.97\linewidth]{ui/#1.png}}}%
    {\fbox{\parbox[c][4cm][c]{0.9\linewidth}{\centering chưa chụp: \detokenize{#1}}}\PackageWarning{solieu}{Thieu so lieu 'ui/#1'}}
  \caption{#2}\label{fig:#3}
  \end{figure}}
% \bang{tệp}{chú thích}{nhãn}: bảng do build.py sinh trong generated/.
\newcommand{\bang}[3]{%
  \begin{table}[H]\centering\caption{#2}\label{tab:#3}
  \IfFileExists{generated/#1.tex}%
    {{\small\setlength{\tabcolsep}{4.5pt}\begin{adjustbox}{max width=\linewidth}\input{generated/#1}\end{adjustbox}}%
     \IfFileExists{generated/#1_note.tex}{\par\smallskip{\footnotesize\input{generated/#1_note}\par}}{}}%
    {\fbox{chưa sinh bảng: \detokenize{#1}}\PackageWarning{solieu}{Thieu so lieu 'tab/#1'}}
  \end{table}}

% Số liệu: \SL{tệp/khoá/.../lá}, sinh tự động từ outputs/metrics/*.json bởi build.py.
\newcommand{\defSL}[2]{\expandafter\def\csname sl@#1\endcsname{#2}}
\newcommand{\SL}[1]{\ifcsname sl@#1\endcsname\csname sl@#1\endcsname\else
  \PackageWarning{solieu}{Thieu so lieu '#1'}\textbf{??}\fi}
\input{generated/so_lieu}

\renewcommand{\arraystretch}{1.16}
\newcolumntype{L}[1]{>{\raggedright\arraybackslash}p{#1}}
\newcolumntype{C}[1]{>{\centering\arraybackslash}p{#1}}
\newcolumntype{R}[1]{>{\raggedleft\arraybackslash}p{#1}}

\begin{document}

% =====================================================================
%  TRANG BÌA (mẫu Assignment 04: khung lùi 6,5 mm, nội dung đặt tuyệt đối trong khung)
% =====================================================================
\begin{titlepage}
\thispagestyle{empty}
\tikzexternaldisable
\begin{tikzpicture}[remember picture,overlay]
  \node[anchor=north west, inner sep=0pt] at ([xshift=6.5mm,yshift=-6.5mm]current page.north west)
    {\includegraphics[width=197mm,height=284mm]{cover-frame.png}};
  \node[anchor=north, inner sep=0pt, text width=161mm, align=center]
    at ([yshift=-(6.5mm+20mm+6mm)]current page.north) {%
    {\bfseries\fontsize{15}{21}\selectfont HỌC VIỆN CÔNG NGHỆ BƯU CHÍNH VIỄN THÔNG\par}
    \vspace{4mm}
    {\bfseries\fontsize{14}{19}\selectfont KHOA CÔNG NGHỆ THÔNG TIN 1\par}
    \vspace{4mm}
    {\bfseries\fontsize{13.5}{18}\selectfont HỌC PHẦN PHÁT TRIỂN CÁC HỆ THỐNG THÔNG MINH\par}
    \vspace{12mm}
    \includegraphics[height=40mm]{ptit-logo.png}\par
    \vspace{10mm}
    {\bfseries\fontsize{22}{27}\selectfont\addfontfeature{LetterSpace=2}BÁO CÁO TIỂU LUẬN MÔN HỌC\par}
    \vspace{5mm}
    {\bfseries\fontsize{13}{19.5}\selectfont ĐỀ TÀI: NGHIÊN CỨU VÀ TRIỂN KHAI CÁC MÔ HÌNH\\
     HỌC MÁY CƠ BẢN, MẠNG NƠ-RON TÍCH CHẬP (CNN)\\
     VÀ MẠNG NƠ-RON HỒI QUY (RNN) TRONG HỆ THỐNG THÔNG MINH\par}
    \vspace{11mm}
    {\fontsize{13}{16}\selectfont
     \renewcommand{\arraystretch}{1.8}
     \begin{tabular}{@{}p{52mm}l@{}}
       \textbf{Lớp:}                  & D23CTPM01 \\
       \textbf{Nhóm:}                 & Nhóm lớp CT -- Nhóm 07 \\
       \textbf{Họ và tên:}            & Nguyễn Duy Nghĩa \\
       \textbf{Mã sinh viên:}         & B23DCCN600 \\
       \textbf{Giảng viên hướng dẫn:} & PGS.TS Trần Đình Quế \\
       \textbf{Học kỳ:}               & Học kỳ 1 năm học 2026 -- 2027 \\
     \end{tabular}\par}};
  \node[anchor=south, inner sep=0pt] at ([yshift=6.5mm+10mm+14mm]current page.south)
    {\bfseries\itshape\fontsize{13}{18}\selectfont Hà Nội, Tháng 10/2026};
\end{tikzpicture}
\tikzexternalenable
\null
\end{titlepage}

% =====================================================================
%  DANH MỤC
% =====================================================================
\pagenumbering{roman}
\setcounter{page}{1}
\tableofcontents
\cleardoublepage
\phantomsection
\addcontentsline{toc}{chapter}{Danh mục Hình vẽ}
\listoffigures
\cleardoublepage
\phantomsection
\addcontentsline{toc}{chapter}{Danh mục Bảng biểu}
\listoftables
\cleardoublepage
\phantomsection
\addcontentsline{toc}{chapter}{Danh mục Đoạn mã nguồn}
\listoflistings

% =====================================================================
%  PHẦN THÂN
% =====================================================================
\cleardoublepage
\pagenumbering{arabic}
\setcounter{page}{1}
% =====================================================================
%  PHẦN MỞ ĐẦU
% =====================================================================
\chapter*{Phần Mở Đầu}
\addcontentsline{toc}{chapter}{Phần Mở Đầu}
\markboth{Phần Mở Đầu}{}

\section*{1. Bối cảnh và tính cấp thiết của đề tài}

Trí tuệ nhân tạo (Artificial Intelligence -- AI) và học máy (Machine Learning -- ML) đã rời phòng thí nghiệm để
trở thành một thành phần thường trực của phần mềm nghiệp vụ: hệ thống sàng lọc bệnh, công cụ định giá tài sản,
bộ lọc phản hồi khách hàng, ứng dụng nhận dạng ảnh, mô hình dự báo tài chính. Điểm chung của các hệ thống này là
chúng không được lập trình bằng luật viết tay mà \emph{học từ dữ liệu}. Sự chuyển dịch đó gắn với ba điều kiện
xuất hiện cùng lúc trong thập niên 2010: dữ liệu lớn, phần cứng tính toán song song (GPU) và các thuật toán tối ưu
cho mạng nhiều tầng \cite{goodfellow2016}.

Một kỹ sư phát triển hệ thống thông minh hiếm khi chỉ gặp một dạng dữ liệu. Trong thực tế có ít nhất ba hình thái,
mỗi hình thái hợp với một họ mô hình riêng:

\begin{itemize}
  \item \textbf{Dữ liệu bảng} (y tế, bất động sản, tài chính doanh nghiệp): mỗi dòng là một đối tượng độc lập với
    các thuộc tính số và phân loại. Đây là địa hạt của các thuật toán học máy cổ điển (hồi quy tuyến tính,
    Logistic Regression, cây quyết định, Random Forest, SVM) và của mạng nơ-ron đa tầng (MLP).
  \item \textbf{Dữ liệu không gian} (ảnh): các điểm ảnh kề nhau tương quan chặt, và cùng một mẫu hình có thể xuất
    hiện ở bất kỳ vị trí nào. Mạng nơ-ron tích chập (CNN) khai thác đúng hai tính chất này nhờ chia sẻ trọng số
    và vùng tiếp nhận cục bộ.
  \item \textbf{Dữ liệu chuỗi} (giá chứng khoán, nhật ký hành vi người dùng): thứ tự thời gian mang thông tin, các
    mẫu không còn độc lập. Mạng nơ-ron hồi quy (RNN) và các biến thể có cổng (LSTM, GRU, BiLSTM) mang một trạng
    thái ẩn đi qua từng bước thời gian để ghi nhớ ngữ cảnh.
\end{itemize}

Từ yêu cầu đó, tiểu luận \emph{“Nghiên cứu và triển khai các mô hình học máy cơ bản, mạng nơ-ron tích chập (CNN)
và mạng nơ-ron hồi quy (RNN) trong hệ thống thông minh”} đặt mục tiêu đi trọn một vòng từ lý thuyết đến sản phẩm
cho cả ba hình thái dữ liệu: hiểu toán học của mô hình, tự cài đặt lại nó, huấn luyện trên dữ liệu thật, đo đạc
khách quan rồi đóng gói thành ứng dụng web chạy được.

\section*{2. Mục tiêu và nhiệm vụ}

Mục tiêu tổng quát là xây dựng nền tảng kiến thức vững cùng năng lực thực hành trên toàn bộ tiến trình từ học máy
cổ điển đến các kiến trúc học sâu xử lý không gian và thời gian. Sáu nhiệm vụ cụ thể:

\begin{description}[leftmargin=2.2em, style=nextline]
  \item[Nhiệm vụ 1 -- Khảo cứu lịch sử AI.] Hệ thống hóa các mốc phát triển từ thập niên 1940 đến nay, phân tích
    nguyên nhân hai “mùa đông AI”, sự chuyển dịch sang học máy thống kê, cuộc cách mạng học sâu và xu hướng hiện tại
    (mô hình nền tảng, RAG, hệ đa tác tử).
  \item[Nhiệm vụ 2 -- Cơ sở lý thuyết.] Trình bày hàm mất mát và cách tối ưu của từng họ mô hình: lan truyền ngược
    trong MLP, phép tích chập và pooling trong CNN, lan truyền ngược qua thời gian và cơ chế cổng trong RNN.
  \item[Nhiệm vụ 3 -- Khảo sát dữ liệu thật.] Phân tích khám phá (EDA) và tiền xử lý \textbf{bảy} tập dữ liệu thuộc
    năm lĩnh vực: y tế (\SL{eda/diabetes/n} hồ sơ tiểu đường), bất động sản Việt Nam (\SL{eda/house/n} tin đăng),
    thương mại điện tử (\SL{eda/reviews/n} nhận xét khách hàng), thị giác máy tính (SVHN
    \SL{eda/svhn/n_train}\,+\,\SL{eda/svhn/n_test} ảnh, EuroSAT \SL{eda/eurosat/n_train}\,+\,\SL{eda/eurosat/n_test} ảnh vệ tinh),
    tài chính (\SL{eda/stock/n_total} phiên giao dịch của bốn mã VN30) và hành vi người dùng
    (\SL{eda/churn/n_users} chuỗi nghe nhạc của KKBox).
  \item[Nhiệm vụ 4 -- Hiện thực trên ba nền tảng.] Cài đặt mô hình bằng \textbf{NumPy thuần} (tự viết lan truyền tiến,
    lan truyền ngược, tích chập, tế bào LSTM), bằng \textbf{Keras/TensorFlow} (API cấp cao) và bằng
    \textbf{PyTorch} (lớp \texttt{nn.Module}, vòng huấn luyện tự viết, chạy GPU).
  \item[Nhiệm vụ 5 -- Đánh giá và đối chuẩn.] So sánh các mô hình trên cùng một cách chia dữ liệu theo độ chính xác
    (Accuracy, Precision, Recall, F1, ROC-AUC, MAE, RMSE, $R^2$), số tham số, dung lượng tệp, thời gian huấn luyện
    và độ trễ suy luận đo trong cùng điều kiện phần cứng.
  \item[Nhiệm vụ 6 -- Triển khai.] Đóng gói các mô hình tốt nhất thành một ứng dụng web Flask gồm sáu phân hệ, nhận
    dữ liệu người dùng và trả kết quả dự báo theo thời gian thực.
\end{description}

\section*{3. Đối tượng và phương pháp nghiên cứu}

\paragraph{Đối tượng.} Các mô hình học máy cơ bản (Linear/Ridge Regression, Logistic Regression, Random Forest, SVM,
MLP); các kiến trúc CNN (Custom CNN, VGG-11, MobileNetV2, ResNet-18 và 1D CNN); các kiến trúc RNN (Simple RNN, LSTM,
GRU, BiLSTM); cùng bảy tập dữ liệu kể trên. Mọi mô hình học sâu được \emph{huấn luyện từ đầu} với trọng số ngẫu
nhiên, không dùng trọng số tiền huấn luyện, để kết quả phản ánh đúng năng lực của từng kiến trúc trên dữ liệu đích.

\paragraph{Phương pháp.} Kết hợp nghiên cứu lý thuyết với kiểm chứng thực nghiệm có kiểm soát:
\begin{itemize}
  \item cùng một hạt giống ngẫu nhiên (\texttt{seed = 42}) và cùng một cách chia huấn luyện/kiểm thử cho mọi mô hình
    trên một tập dữ liệu; bộ tiền xử lý chỉ được khớp trên phần huấn luyện;
  \item chia phân tầng theo nhãn cho bài toán phân loại mất cân bằng, chia theo trục thời gian (không xáo trộn) cho
    chuỗi giá chứng khoán để tránh rò rỉ thông tin tương lai;
  \item tập kiểm định tách riêng cho dừng sớm và chọn ngưỡng; tập kiểm thử chỉ dùng một lần để báo cáo;
  \item mỗi kết luận so sánh đi kèm một mốc tham chiếu (ví dụ dự báo “giữ nguyên giá hôm nay” cho chứng khoán);
  \item mọi con số trong báo cáo được sinh tự động từ tệp kết quả JSON của chính lần chạy, không chép tay.
\end{itemize}

\paragraph{Phân bổ phần cứng.} Máy thực nghiệm có CPU 32 luồng và GPU NVIDIA GeForce RTX 4060 Laptop (8 GB). PyTorch
dùng GPU; TensorFlow~2.21 trên Windows không hỗ trợ GPU nên các mô hình Keras chạy CPU. Các tiến trình độc lập chạy
song song: scikit-learn, mạng NumPy và Keras trên CPU (giới hạn số luồng mỗi tiến trình), PyTorch trên GPU.

\section*{4. Bố cục báo cáo}

\begin{itemize}
  \item \textbf{Chương 1 -- Lịch sử phát triển trí tuệ nhân tạo:} bức tranh hơn 80 năm từ nơ-ron hình thức đến hệ đa tác tử.
  \item \textbf{Chương 2 -- Các kỹ thuật học máy cơ bản:} lý thuyết bốn họ thuật toán, EDA ba tập dữ liệu bảng và văn bản,
    pipeline scikit-learn, MLP trên ba nền tảng, đối chuẩn và Web App.
  \item \textbf{Chương 3 -- Mạng nơ-ron tích chập:} phép tích chập, vùng cảm nhận, pooling, ba thế hệ kiến trúc;
    SVHN, EuroSAT và 1D CNN; hiện thực bằng NumPy, Keras và PyTorch; đối chuẩn độ chính xác -- chi phí; Web App thị giác.
  \item \textbf{Chương 4 -- Mạng nơ-ron hồi quy:} RNN, BPTT, triệt tiêu đạo hàm, cắt tỉa gradient, LSTM, GRU, BiLSTM;
    chứng khoán VN30 và KKBox; hiện thực ba nền tảng; đối chuẩn PyTorch -- Keras; Web App dự báo.
\end{itemize}
Phần Kết luận tổng hợp các phát hiện và hướng phát triển. Mã nguồn đầy đủ nằm trong thư mục
\texttt{src/TieuLuan 01} của kho lưu trữ; mỗi đoạn mã trong báo cáo được trích nguyên văn từ đó.

% =====================================================================
%  CHƯƠNG 1 — LỊCH SỬ PHÁT TRIỂN TRÍ TUỆ NHÂN TẠO
% =====================================================================
\chapter{Lịch Sử Phát Triển Trí Tuệ Nhân Tạo}
\label{ch:lich-su}

\section{Tổng quan bức tranh hơn tám thập kỷ của AI}

Lịch sử trí tuệ nhân tạo không phải một đường đi lên đều đặn. Đó là chuỗi các chu kỳ kỳ vọng: mỗi đột phá kéo theo
lời hứa vượt quá năng lực thật của công nghệ, rồi một giai đoạn thoái trào khi tiền tài trợ cạn, gọi là
\emph{mùa đông AI}. Qua các chu kỳ đó, cách máy tính “suy nghĩ” đổi ba lần: từ suy diễn trên ký hiệu do con người
viết, sang mô hình xác suất học từ dữ liệu, rồi tới học biểu diễn nhiều tầng của học sâu. Hình~\ref{fig:timeline}
đặt mười sáu mốc then chốt lên cùng một trục thời gian, kèm các kỷ nguyên tương ứng.

\hinhtikz{f1_timeline}{Dòng thời gian tiến hóa của trí tuệ nhân tạo (1943 -- 2026): các mốc then chốt và các kỷ
nguyên công nghệ. Hai dải xám là hai mùa đông AI (1974 -- 1980 và 1987 -- 1993), dải cam giữa chúng là thời kỳ hệ chuyên
gia (1980 -- 1987).}{timeline}

Hiểu lịch sử không chỉ để biết ai làm gì, năm nào. Mỗi mô hình được hiện thực ở các chương sau đều là lời đáp cho
một thất bại cụ thể trong quá khứ: MLP và lan truyền ngược trả lời giới hạn XOR của Perceptron, SVM và Random Forest
trả lời sự giòn của hệ chuyên gia, LSTM trả lời hiện tượng triệt tiêu đạo hàm của RNN, ResNet trả lời sự thoái hóa của
mạng rất sâu.

\section{Khởi nguồn và kỷ nguyên logic biểu tượng (1940s -- 1960s)}

\subsection{Ý niệm về máy tính có tư duy và phép thử Turing}

Năm 1950, Alan Turing công bố bài báo \emph{“Computing Machinery and Intelligence”} trên tạp chí \emph{Mind}
\cite{turing1950}, mở đầu bằng câu hỏi “Máy móc có thể suy nghĩ không?”. Để tránh tranh cãi triết học về khái niệm
“ý thức”, Turing thay câu hỏi bằng một thí nghiệm hành vi, \emph{trò chơi bắt chước}: một người thẩm vấn trò chuyện
bằng văn bản với hai đối tượng giấu mặt, một người và một máy. Nếu người thẩm vấn không phân biệt được đâu là máy,
ta nói máy đã thể hiện hành vi thông minh. Phép thử Turing là phương pháp luận đầu tiên định lượng hóa trí tuệ nhân
tạo bằng hành vi quan sát được, và đến nay vẫn là điểm quy chiếu khi bàn về các mô hình ngôn ngữ lớn.

\subsection{Mô hình nơ-ron hình thức và Hội thảo Dartmouth (1956)}

Trước cả Turing, năm 1943, nhà thần kinh học Warren McCulloch và nhà logic học Walter Pitts đề xuất mô hình
nơ-ron hình thức \cite{mcculloch1943}: một nơ-ron tính tổng có trọng số các tín hiệu nhị phân rồi so với một ngưỡng,
\begin{equation}
  y = f\Big(\sum_{i=1}^{n} w_i x_i - \theta\Big), \qquad x_i \in \{0, 1\},
  \label{eq:mcp}
\end{equation}
với $f$ là hàm bước Heaviside. Hai tác giả chứng minh mạng các nơ-ron như vậy thực hiện được mọi hàm logic AND, OR,
NOT. Năm 1949, Donald Hebb bổ sung quy tắc học “các nơ-ron cùng phát xung thì liên kết mạnh lên”, cơ sở sinh học đầu
tiên cho việc \emph{thay đổi} trọng số.

Ngành AI chính thức ra đời mùa hè năm 1956 tại Hội thảo Dartmouth do John McCarthy, Marvin Minsky, Claude Shannon và
Nathaniel Rochester tổ chức. McCarthy đặt tên “Artificial Intelligence” với giả thuyết: mọi khía cạnh của việc học hay
của trí tuệ, về nguyên tắc, đều mô tả được chính xác tới mức một cỗ máy có thể mô phỏng nó.

\subsection{Perceptron của Rosenblatt và các hệ suy diễn ký hiệu đầu tiên}

Năm 1958, Frank Rosenblatt công bố Perceptron \cite{rosenblatt1958}, mô hình nơ-ron đầu tiên \emph{tự học} trọng số
từ dữ liệu có nhãn theo quy tắc
\begin{equation}
  \mathbf{w}^{(t+1)} = \mathbf{w}^{(t)} + \eta\,(y - \hat y)\,\mathbf{x},
  \label{eq:perceptron}
\end{equation}
với $\eta > 0$ là tốc độ học. Định lý hội tụ Perceptron bảo đảm thuật toán dừng sau hữu hạn bước nếu dữ liệu khả
tách tuyến tính. Song song với hướng kết nối (connectionism) này, hướng biểu tượng (symbolic AI) cũng gặt hái thành
công: chương trình Logic Theorist (1956) của Newell, Simon và Shaw chứng minh được 38 trên 52 định lý ở chương 2 của
\emph{Principia Mathematica}; General Problem Solver (1959) giải các bài toán đố bằng phân tích phương tiện -- mục đích.

\section{“Mùa đông AI” lần thứ nhất (1970s)}

\subsection{Giới hạn toán học của Perceptron và cuốn sách của Minsky \& Papert (1969)}

Năm 1969, Minsky và Papert xuất bản \emph{Perceptrons} \cite{minsky1969}, chứng minh chặt chẽ rằng Perceptron một tầng
chỉ học được các hàm khả tách tuyến tính, và vì vậy thất bại với cả hàm XOR đơn giản (Bảng~\ref{tab:xor}).

\begin{table}[H]\centering
\caption{Bảng chân lý của hàm XOR: hai lớp nằm trên hai đường chéo nên không đường thẳng nào tách được.}\label{tab:xor}
\begin{tabular}{cccl}
\toprule
$x_1$ & $x_2$ & $y = x_1 \oplus x_2$ & Vị trí trên mặt phẳng $\mathbb{R}^2$ \\
\midrule
0 & 0 & 0 & gốc tọa độ $(0, 0)$ -- nhãn âm \\
0 & 1 & 1 & điểm $(0, 1)$ -- nhãn dương \\
1 & 0 & 1 & điểm $(1, 0)$ -- nhãn dương \\
1 & 1 & 0 & điểm $(1, 1)$ -- nhãn âm \\
\bottomrule
\end{tabular}
\end{table}

Muốn tách $\{(0,1), (1,0)\}$ khỏi $\{(0,0), (1,1)\}$ cần một đường thẳng $w_1x_1 + w_2x_2 + b = 0$ để hai nhóm nằm hai
phía, nhưng từ $b < 0$, $w_1 + b > 0$, $w_2 + b > 0$ suy ra $w_1 + w_2 + b > -b > 0$, mâu thuẫn với yêu cầu điểm
$(1,1)$ phải âm. Minsky và Papert biết mạng nhiều tầng giải được XOR, nhưng thời đó chưa ai có cách huấn luyện các
tầng ẩn. Cuốn sách khiến nguồn tài trợ cho hướng nơ-ron gần như đóng băng suốt hơn một thập kỷ.

\subsection{Báo cáo Lighthill (1973) và sự cắt giảm tài trợ của DARPA}

Năm 1973, Báo cáo Lighthill \cite{lighthill1973} gửi Hội đồng Nghiên cứu Khoa học Anh kết luận rằng các hệ AI ký hiệu
chỉ chạy được trong “thế giới đồ chơi” và bất lực trước bài toán thật vì \emph{bùng nổ tổ hợp}: số trạng thái cần tìm
kiếm tăng theo hàm mũ của kích thước bài toán. Chính phủ Anh cắt tài trợ AI ở hầu hết các trường đại học; cùng lúc, tu
chính án Mansfield khiến DARPA của Hoa Kỳ ngừng chi cho nghiên cứu không có ứng dụng quân sự trực tiếp. Giai đoạn
1974 -- 1980 được gọi là mùa đông AI thứ nhất.

\section{Sự trỗi dậy của hệ chuyên gia và “mùa đông AI” lần thứ hai (1980s -- 1990s)}

\subsection{Kỷ nguyên vàng của hệ chuyên gia}

Đầu thập niên 1980, AI hồi sinh khi chuyển từ tìm kiếm tổng quát sang \emph{tri thức chuyên ngành hẹp}. Một hệ chuyên
gia tách hai thành phần: \textbf{cơ sở tri thức} gồm sự kiện và luật dạng \texttt{IF} điều kiện \texttt{THEN} kết luận,
và \textbf{máy suy diễn} chạy luật theo chiều tiến (từ dữ kiện ra kết luận) hoặc chiều lùi (từ giả thuyết truy về dữ
kiện). Các hệ tiêu biểu gồm DENDRAL (xác định cấu trúc phân tử từ phổ khối), MYCIN (chẩn đoán nhiễm trùng máu, độ chính
xác ngang bác sĩ chuyên khoa) và XCON/R1 của Digital Equipment Corporation, tự cấu hình máy VAX và được ước tính tiết
kiệm khoảng 40 triệu USD mỗi năm. Thành công của XCON kéo theo làn sóng đầu tư vào máy Lisp chuyên dụng và dự án Máy
tính Thế hệ Thứ năm của Nhật Bản.

\subsection{Sự tái sinh của mạng nơ-ron: thuật toán lan truyền ngược (1986)}

Năm 1986, Rumelhart, Hinton và Williams công bố trên \emph{Nature} thuật toán lan truyền ngược cho mạng nhiều tầng
\cite{rumelhart1986}. Dùng quy tắc dây chuyền với hàm kích hoạt khả vi như sigmoid $\sigma(z) = 1/(1+e^{-z})$, đạo hàm
của mất mát $L$ theo một trọng số bất kỳ ở tầng ẩn $l$ tách thành tích của sai số truyền ngược và kích hoạt tầng trước:
\begin{equation}
  \frac{\partial L}{\partial w^{(l)}_{ij}} = \delta^{(l)}_i \, a^{(l-1)}_j .
  \label{eq:bp}
\end{equation}
Kết quả này phá “lời nguyền XOR”: mạng hai tầng học được biểu diễn ẩn phi tuyến. Cùng giai đoạn, Fukushima
(1980) đề xuất Neocognitron \cite{fukushima1980} với các tầng trích đặc trưng cục bộ (S-cells) xen kẽ các tầng bất biến
vị trí (C-cells), tiền thân trực tiếp của CNN. Chương~\ref{ch:hoc-may} hiện thực lại công thức~\eqref{eq:bp} bằng NumPy.

\subsection{“Mùa đông AI” lần thứ hai (1987 -- 1993)}

Hệ chuyên gia sớm lộ ba điểm yếu cấu trúc. \emph{Nút thắt thu nhận tri thức}: phỏng vấn chuyên gia để mã hóa hàng vạn
luật vừa đắt vừa chậm, lại sinh mâu thuẫn. \emph{Tính giòn}: hệ thống trả lời sai hoặc lặp vô tận khi gặp tình huống
nằm ngoài các luật đã viết, chi phí bảo trì vượt lợi ích. \emph{Sụp đổ phần cứng}: máy trạm Unix và vi xử lý x86 rẻ
hơn khiến máy Lisp lỗi thời, kéo theo phá sản dây chuyền. Hình~\ref{fig:hype} minh họa định tính các chu kỳ kỳ vọng và
hai mùa đông AI.

\hinhtikz{f1_hype}{Chu kỳ kỳ vọng và hai thời kỳ “mùa đông AI” (minh họa định tính, không phải số đo).}{hype}

Sau mỗi lần thoái trào, cộng đồng tái cấu trúc phương pháp luận theo hướng thực dụng hơn và có cơ sở toán học chắc hơn.

\section{Kỷ nguyên học máy thống kê (1990s -- 2010)}

Bài học rút ra từ thất bại của hệ chuyên gia: thay vì cố viết tay mọi quy tắc của thế giới, hãy để máy học quy luật
xác suất trực tiếp từ dữ liệu.

\subsection{Lý thuyết học thống kê và máy vector hỗ trợ (SVM)}

Vapnik và Chervonenkis xây dựng lý thuyết học thống kê với nguyên lý cực tiểu hóa rủi ro cấu trúc: kiểm soát cùng lúc
sai số huấn luyện và độ phức tạp của lớp hàm (đo bằng chiều VC). Đỉnh cao là máy vector hỗ trợ \cite{cortes1995}, tìm
siêu phẳng có \emph{lề lớn nhất} tới các điểm gần nhất:
\begin{equation}
  \min_{\mathbf{w}, b}\ \tfrac{1}{2}\lVert \mathbf{w} \rVert^2 \quad \text{với}\quad y_i(\mathbf{w}^\top\mathbf{x}_i + b) \ge 1,\ \forall i.
\end{equation}
Nhờ thủ thuật hàm nhân, ví dụ nhân RBF $K(\mathbf{x}, \mathbf{x}') = \exp(-\gamma\lVert\mathbf{x}-\mathbf{x}'\rVert^2)$,
SVM tách được dữ liệu phi tuyến bằng cách ngầm ánh xạ lên không gian vô hạn chiều. Bài toán tối ưu là lồi nên không
mắc cực tiểu địa phương; SVM thống trị nhận dạng ảnh và phân loại văn bản suốt hơn một thập kỷ.

\subsection{Cây quyết định và các phương pháp học tập hợp}

Cây quyết định (CART, C4.5) chia không gian bằng chuỗi câu hỏi đúng/sai, dễ diễn giải. Hai cách kết hợp nhiều mô hình
yếu thành một mô hình mạnh xác lập vị thế trên dữ liệu bảng: \textbf{Boosting} (AdaBoost \cite{freund1997}, sau đó
Gradient Boosting) huấn luyện tuần tự các mô hình sửa lỗi của nhau, còn \textbf{Random Forest} \cite{breiman2001}
lấy trung bình hàng trăm cây học trên các mẫu bootstrap với tập con đặc trưng ngẫu nhiên, giảm phương sai mà không
tăng độ chệch. Chương~\ref{ch:hoc-may} cho thấy Random Forest đến nay vẫn là đối thủ khó vượt trên dữ liệu bảng.

\subsection{LeNet-5 (1998): nền tảng của mạng tích chập hiện đại}

Năm 1998, LeCun và cộng sự công bố LeNet-5 \cite{lecun1998}, ghép ba ý tưởng: vùng tiếp nhận cục bộ, chia sẻ trọng số
và lấy mẫu giảm (pooling). LeNet-5 đọc hàng triệu tấm séc ngân hàng mỗi ngày ở Hoa Kỳ, nhưng giới hạn CPU và thiếu dữ
liệu ảnh lớn khiến nó chưa mở rộng được sang ảnh tự nhiên phức tạp.

\section{Cuộc cách mạng học sâu (2012 -- 2020)}

\subsection{Bước ngoặt ImageNet 2012 và AlexNet}

Tại cuộc thi ILSVRC 2012, AlexNet của Krizhevsky, Sutskever và Hinton \cite{krizhevsky2012} đạt sai số top-5 15,3\,\%,
bỏ xa vị trí thứ hai (26,2\,\%) dùng đặc trưng thủ công SIFT/HOG cùng SVM. Ba yếu tố hội tụ: \textbf{dữ liệu lớn}
(ImageNet hơn 1,2 triệu ảnh, 1.000 lớp), \textbf{GPU} với CUDA tăng tốc nhân ma trận hàng chục lần so với CPU, và
\textbf{các cải tiến thuật toán} như hàm kích hoạt ReLU $f(x) = \max(0, x)$ không bão hòa ở vùng dương và điều chuẩn
Dropout \cite{srivastava2014} ngắt ngẫu nhiên nơ-ron để chống đồng thích nghi.

\subsection{Sự tiến hóa của các kiến trúc thị giác: VGG, ResNet, MobileNet}

\textbf{VGGNet} \cite{simonyan2015} chuẩn hóa thiết kế bằng cách xếp chồng các bộ lọc $3{\times}3$: hai tầng $3{\times}3$
có cùng vùng nhìn với một tầng $5{\times}5$ nhưng ít hơn 28\,\% trọng số và thêm một hàm phi tuyến. \textbf{ResNet}
\cite{he2016} giải quyết hiện tượng thoái hóa (mạng sâu hơn lại có sai số huấn luyện cao hơn) bằng kết nối tắt
$\mathbf{y} = F(\mathbf{x}) + \mathbf{x}$, cho phép huấn luyện mạng 152 tầng và đưa sai số top-5 ImageNet xuống 3,57\,\%.
\textbf{MobileNet} \cite{howard2017,sandler2018} tách tích chập chuẩn thành depthwise và pointwise, giảm chi phí tính
toán 8 -- 9 lần để chạy trên điện thoại. Cả ba kiến trúc được hiện thực và đối chuẩn ở Chương~\ref{ch:cnn}.

\subsection{Mô hình hóa chuỗi: từ RNN, LSTM đến Transformer}

RNN duy trì một trạng thái ẩn qua các bước thời gian, nhưng lan truyền ngược qua thời gian nhân liên tiếp các ma trận
Jacobian nên gradient suy giảm theo hàm mũ \cite{bengio1994}. LSTM \cite{hochreiter1997} thêm ô nhớ được điều khiển bởi
ba cổng (quên, vào, ra), tạo đường đi cộng cho gradient; GRU \cite{cho2014} rút gọn còn hai cổng. Năm 2017, Transformer
\cite{vaswani2017} bỏ hẳn liên kết hồi quy, dùng cơ chế tự chú ý đa đầu để mọi vị trí trong chuỗi nhìn thấy nhau trực
tiếp và tính song song hoàn toàn trên GPU, mở đường cho các mô hình ngôn ngữ khổng lồ. Chương~\ref{ch:rnn} hiện thực
và so sánh bốn kiến trúc hồi quy.

\section{Kỷ nguyên AI tạo sinh, mô hình nền tảng và hệ đa tác tử (2020 -- nay)}

\subsection{Mô hình ngôn ngữ lớn và định luật tỷ lệ}

Kaplan và cộng sự \cite{kaplan2020}, rồi Hoffmann và cộng sự (Chinchilla) \cite{hoffmann2022} cho thấy mất mát của mô
hình ngôn ngữ giảm theo luật lũy thừa khá đều với ba đại lượng: số tham số, lượng dữ liệu và ngân sách tính toán. Ở quy
mô hàng trăm tỷ tham số xuất hiện các năng lực nổi trội: học trong ngữ cảnh từ vài ví dụ mà không cập nhật trọng số,
suy luận từng bước (chain-of-thought), sinh mã nguồn. Hình~\ref{fig:nested} tóm tắt quan hệ bao hàm giữa các thế hệ.

\hinhtikz{f1_nested}{Quan hệ bao hàm giữa trí tuệ nhân tạo, học máy, học sâu và AI tạo sinh / hệ đa tác tử.}{nested}

\subsection{Giới hạn của mô hình thuần tham số và vai trò của đồ thị tri thức, RAG}

Mô hình ngôn ngữ chỉ lưu tri thức trong trọng số nên có ba giới hạn: \emph{ảo giác} (sinh câu trả lời trôi chảy nhưng
sai sự thật, vì mô hình tối ưu xác suất từ tiếp theo chứ không kiểm chứng), \emph{tri thức đóng băng} tại thời điểm
tiền huấn luyện, và \emph{khó truy vết nguồn} trong các lĩnh vực cần trích dẫn như y tế, pháp luật. Xu hướng khắc phục
là ghép thêm bộ nhớ phi tham số: \textbf{đồ thị tri thức} lưu sự kiện dưới dạng bộ ba (thực thể, quan hệ, thực thể) để
suy diễn tường minh; \textbf{mạng nơ-ron đồ thị} mở rộng học sâu sang dữ liệu có cấu trúc đồ thị bằng truyền thông điệp
giữa các nút; \textbf{RAG} \cite{lewis2020} truy xuất tài liệu liên quan rồi đưa vào ngữ cảnh trước khi sinh câu trả lời.

\subsection{Xu hướng: hệ đa tác tử tự trị}

Bước tiếp theo là các \emph{tác tử} (agent) tự chủ có chu trình nhận thức đầy đủ: quan sát, lập kế hoạch và suy luận,
gọi công cụ hay API, rồi tự đánh giá và sửa sai. Nhiều tác tử chuyên biệt (lập kế hoạch, nghiên cứu, lập trình, kiểm thử)
phối hợp và phản biện nhau trong một hệ đa tác tử giải được những việc kỹ thuật mà một mô hình đơn lẻ không làm nổi.
Đây là định hướng then chốt của hệ thống thông minh thế hệ mới, và cũng là nội dung các chương sau của học phần
(đồ thị tri thức, mạng nơ-ron đồ thị, hệ đa tác tử) sẽ kế thừa từ nền tảng mô hình hóa được xây dựng trong tiểu luận này.

% =====================================================================
%  CHƯƠNG 2 — CÁC KỸ THUẬT HỌC MÁY CƠ BẢN
% =====================================================================
\chapter{Các Kỹ Thuật Học Máy Cơ Bản}
\label{ch:hoc-may}

\section{Cơ sở lý thuyết các thuật toán học máy cốt lõi}

Học máy có giám sát giải hai lớp bài toán: \emph{hồi quy} (dự báo một giá trị liên tục) và \emph{phân loại} (gán
nhãn rời rạc). Cho tập huấn luyện $\{(\mathbf{x}^{(i)}, y^{(i)})\}_{i=1}^{N}$ với $\mathbf{x} \in \mathbb{R}^{D}$, ta chọn
một họ hàm $f_\theta$ và tìm tham số $\theta$ cực tiểu hóa một hàm mất mát trung bình trên dữ liệu. Bốn họ thuật toán
dưới đây khác nhau ở họ hàm và cách tối ưu \cite{raschka2015,goodfellow2016}.

\subsection{Hồi quy tuyến tính (Linear Regression)}

Mô hình giả định đầu ra là tổ hợp tuyến tính của đặc trưng:
\begin{equation}
  \hat y = f(\mathbf{x}) = \mathbf{w}^\top\mathbf{x} + b = \sum_{j=1}^{D} w_j x_j + b ,
\end{equation}
và cực tiểu hóa sai số toàn phương trung bình
\begin{equation}
  \mathcal{L}_{\mathrm{MSE}}(\mathbf{w}, b) = \frac{1}{2N}\sum_{i=1}^{N}\big(\hat y^{(i)} - y^{(i)}\big)^2
    = \frac{1}{2N}\lVert X\mathbf{w} + b\mathbf{1} - \mathbf{y}\rVert_2^2 .
\end{equation}
Nghiệm đóng là phương trình chuẩn $\mathbf{w}^* = (X^\top X)^{-1}X^\top\mathbf{y}$; khi $D$ lớn có thể dùng hạ gradient
$\mathbf{w} \leftarrow \mathbf{w} - \eta\,\tfrac{1}{N}X^\top(X\mathbf{w} - \mathbf{y})$. Khi các đặc trưng tương quan mạnh
(đa cộng tuyến), $X^\top X$ gần suy biến và hệ số dao động lớn; \textbf{Ridge} thêm phạt $\alpha\lVert\mathbf{w}\rVert_2^2$
để nghiệm thành $(X^\top X + \alpha I)^{-1}X^\top\mathbf{y}$, luôn khả nghịch.

\subsection{Hồi quy Logistic (Logistic Regression)}

Với nhãn $y \in \{0, 1\}$, tổ hợp tuyến tính được ép về khoảng $(0, 1)$ bằng hàm sigmoid:
\begin{equation}
  p = P(y = 1 \mid \mathbf{x}) = \sigma(z) = \frac{1}{1 + e^{-z}}, \qquad z = \mathbf{w}^\top\mathbf{x} + b .
\end{equation}
Ước lượng hợp lý cực đại dẫn tới hàm mất mát entropy chéo nhị phân
\begin{equation}
  \mathcal{L}_{\mathrm{BCE}} = -\frac{1}{N}\sum_{i=1}^{N}\Big[y^{(i)}\log\hat y^{(i)} + (1 - y^{(i)})\log(1 - \hat y^{(i)})\Big],
  \label{eq:bce}
\end{equation}
có đạo hàm theo $z$ rút gọn thành $\partial\mathcal{L}/\partial z = \hat y - y$, nên
$\nabla_{\mathbf{w}}\mathcal{L} = \tfrac{1}{N}\sum_i (\hat y^{(i)} - y^{(i)})\mathbf{x}^{(i)}$. Mất mát là hàm lồi, mọi
cực tiểu đều là toàn cục. Với $K$ lớp, sigmoid được thay bằng softmax $\hat y_k = e^{z_k}/\sum_j e^{z_j}$ và mất mát là
entropy chéo $-\log \hat y_{y}$. Khi dữ liệu mất cân bằng, trọng số lớp \texttt{class\_weight="balanced"} nhân mất mát của
mỗi mẫu với $N/(K\,N_k)$ để lớp hiếm có tiếng nói ngang lớp phổ biến.

\subsection{Cây quyết định và Rừng ngẫu nhiên (Random Forest)}

Cây quyết định chia không gian bằng chuỗi câu hỏi nhị phân “$x_j \le \theta$?”. Độ thuần của một nút $t$ có các lớp với tỉ
lệ $p_k$ được đo bằng chỉ số Gini hoặc entropy:
\begin{equation}
  \mathrm{Gini}(t) = 1 - \sum_{k=1}^{K} p_k^2, \qquad \mathrm{Entropy}(t) = -\sum_{k=1}^{K} p_k\log_2 p_k .
\end{equation}
Tại mỗi nút, thuật toán chọn cặp $(j, \theta)$ làm giảm tạp chất nhiều nhất:
\begin{equation}
  \Delta\mathrm{Gini} = \mathrm{Gini}(t) - \Big(\frac{N_L}{N_t}\mathrm{Gini}(t_L) + \frac{N_R}{N_t}\mathrm{Gini}(t_R)\Big).
\end{equation}
Với hồi quy, tạp chất là phương sai của nhãn trong nút. Một cây sâu khớp rất sát dữ liệu huấn luyện nhưng có phương sai
lớn. \textbf{Random Forest} \cite{breiman2001} trung bình hóa $B$ cây với hai nguồn ngẫu nhiên: mỗi cây học trên một mẫu
bootstrap (rút có hoàn lại $N$ mẫu) và mỗi lần chia chỉ xét một tập con ngẫu nhiên $\sqrt{D}$ đặc trưng. Hai cơ chế này
làm các cây ít tương quan với nhau, nên trung bình của chúng có phương sai nhỏ hơn hẳn từng cây mà độ chệch gần như giữ
nguyên.

\subsection{Máy vector hỗ trợ (Support Vector Machine)}

SVM \cite{cortes1995} tìm siêu phẳng phân tách có lề lớn nhất, cho phép vi phạm qua biến nới lỏng $\xi_i$:
\begin{equation}
  \min_{\mathbf{w}, b, \boldsymbol{\xi}}\ \frac{1}{2}\lVert\mathbf{w}\rVert^2 + C\sum_{i=1}^{N}\xi_i
  \quad\text{với}\quad y^{(i)}(\mathbf{w}^\top\mathbf{x}^{(i)} + b) \ge 1 - \xi_i,\ \ \xi_i \ge 0 ,
\end{equation}
trong đó $C$ cân bằng giữa lề rộng và lỗi huấn luyện. Bài toán đối ngẫu chỉ phụ thuộc vào tích vô hướng giữa các mẫu, nên
có thể thay bằng một hàm nhân. Nhân RBF
\begin{equation}
  K(\mathbf{x}, \mathbf{x}') = \exp\!\big(-\gamma\lVert\mathbf{x} - \mathbf{x}'\rVert^2\big)
\end{equation}
tương ứng với ánh xạ lên không gian vô hạn chiều. Cái giá là bộ giải phải làm việc với ma trận nhân $N \times N$: thời gian
huấn luyện tăng cỡ $O(N^2)$ tới $O(N^3)$ và mô hình phải lưu toàn bộ vector hỗ trợ. \textbf{SVR} là phiên bản hồi quy, phạt
các điểm nằm ngoài ống sai số rộng $\varepsilon$. Với văn bản TF-IDF hàng nghìn chiều, dữ liệu thường đã gần khả tách tuyến
tính nên \textbf{LinearSVC} (nhân tuyến tính, độ phức tạp tuyến tính theo $N$) là lựa chọn hợp lý hơn.

% ---------------------------------------------------------------------
\section{Khảo sát thực nghiệm và tiền xử lý ba bộ dữ liệu}

Ba tập dữ liệu được chọn để phủ ba dạng bài toán: phân loại nhị phân mất cân bằng (y tế), hồi quy (bất động sản Việt Nam)
và phân loại văn bản ba lớp (thương mại điện tử). Mọi tập đều chia 80\,\% huấn luyện / 20\,\% kiểm thử với
\texttt{random\_state = 42} (phân tầng theo nhãn với bài toán phân loại); mạng MLP tách thêm 10\,\% của phần huấn luyện làm
tập kiểm định cho dừng sớm.

\subsection{Tập dữ liệu 1: Sàng lọc nguy cơ tiểu đường (\SL{eda/diabetes/n} mẫu)}

\begin{itemize}
  \item \textbf{Nguồn:} Diabetes Prediction Dataset trên Kaggle \cite{kaggle_diabetes}, tổng hợp hồ sơ khám bệnh.
  \item \textbf{Quy mô:} \SL{eda/diabetes/n} hồ sơ, \SL{eda/diabetes/n_features} thuộc tính: giới tính, tuổi, tăng huyết áp,
    bệnh tim, tiền sử hút thuốc, chỉ số khối cơ thể BMI, HbA1c (\%), đường huyết (mg/dL); nhãn \texttt{diabetes}
    (0: không mắc, 1: mắc).
\end{itemize}

\bang{tab_sample_diab}{Năm bản ghi ngẫu nhiên của tập dữ liệu tiểu đường.}{sample-diab}

\hinhpdf{f2_diab_target}{0.62}{Phân bố nhãn của tập tiểu đường: lớp dương chỉ chiếm \SL{eda/diabetes/pos_ratePct}\,\%.}{diab-target}
\hinhpdf{f2_diab_kde}{1}{Mật độ xác suất của đường huyết, HbA1c, BMI và tuổi theo tình trạng bệnh.}{diab-kde}
\hinhpdf{f2_diab_corr}{0.72}{Ma trận tương quan Pearson giữa các chỉ số lâm sàng và nhãn tiểu đường.}{diab-corr}

\paragraph{Nhận xét EDA.}
\begin{itemize}
  \item \textbf{Mất cân bằng lớp:} chỉ \SL{eda/diabetes/n_pos} trên \SL{eda/diabetes/n} hồ sơ là dương tính
    (Hình~\ref{fig:diab-target}). Một mô hình ngây thơ luôn đoán “không mắc” đã đạt Accuracy
    $100 - \SL{eda/diabetes/pos_ratePct} \approx 91{,}5\,\%$ mà không phát hiện được ca bệnh nào; vì vậy Accuracy không
    đủ để đánh giá, cần Precision, Recall và F1 của lớp dương cùng ROC-AUC.
  \item \textbf{Hai đặc trưng quyết định:} đường huyết và HbA1c có tương quan với nhãn lần lượt
    $r = \SL{eda/diabetes/corr_glucose}$ và $r = \SL{eda/diabetes/corr_hba1c}$, cao nhất trong các chỉ số
    (Hình~\ref{fig:diab-corr}). Phân bố của người mắc bệnh dịch hẳn sang phải (Hình~\ref{fig:diab-kde}): tỉ lệ mắc bệnh là
    \SL{eda/diabetes/rate_glu200Pct}\,\% khi đường huyết $\ge 200$\,mg/dL, \SL{eda/diabetes/rate_hba65Pct}\,\% khi
    HbA1c $\ge 6{,}6$\,\%, nhưng chỉ \SL{eda/diabetes/rate_hba_lowPct}\,\% khi HbA1c $< 5{,}7$\,\%.
  \item \textbf{Quan hệ dạng ngưỡng:} cả hai chỉ số đều có “ngưỡng gãy” rõ (đặc biệt HbA1c), một dạng phi tuyến mà
    mô hình tuyến tính không biểu diễn trọn được nhưng cây quyết định và mạng nơ-ron thì có thể.
  \item \textbf{Tiền xử lý:} chuẩn hóa z-score bốn biến số, giữ nguyên hai biến nhị phân, mã hóa one-hot giới tính và tiền
    sử hút thuốc (\SL{eda/diabetes/noinfo_smokePct}\,\% hồ sơ ghi “No Info”, được giữ thành một nhóm riêng).
\end{itemize}

\subsection{Tập dữ liệu 2: Định giá bất động sản Việt Nam (\SL{eda/house/n} mẫu)}

\begin{itemize}
  \item \textbf{Nguồn:} Vietnam Housing Dataset \cite{kaggle_vnhouse}, tin đăng bán nhà trên các trang bất động sản Việt Nam.
  \item \textbf{Quy mô:} \SL{eda/house/n} tin đăng thuộc \SL{eda/house/n_province} tỉnh/thành và \SL{eda/house/n_district}
    quận/huyện; các thuộc tính diện tích, mặt tiền, độ rộng đường vào, số tầng, số phòng ngủ, số phòng tắm, hướng nhà,
    hướng ban công, tình trạng pháp lý, nội thất, địa chỉ; nhãn là giá bán (tỷ đồng).
\end{itemize}

\bang{tab_sample_house}{Năm tin đăng ngẫu nhiên của tập dữ liệu giá nhà.}{sample-house}

\hinhpdf{f2_house_dist}{1}{Phân bố giá bán, diện tích gốc và diện tích sau biến đổi $\log(1 + x)$.}{house-dist}
\hinhpdf{f2_house_scatter}{1}{Trái: giá theo diện tích (thang log), tô màu theo độ rộng đường vào. Phải: giá theo sáu
tỉnh/thành có nhiều tin nhất.}{house-scatter}
\hinhpdf{f2_house_corr}{0.62}{Ma trận tương quan giữa các thuộc tính vật lý và giá bán.}{house-corr}

\paragraph{Nhận xét EDA.}
\begin{itemize}
  \item \textbf{Giá đã bị chặn:} giá nằm trong $[\SL{eda/house/price_min}; \SL{eda/house/price_max}]$ tỷ đồng, trung bình
    \SL{eda/house/price_mean}, độ lệch $\SL{eda/house/price_skew}$ (Hình~\ref{fig:house-dist}). Phân bố gần đối xứng nên
    \emph{không} cần biến đổi log cho nhãn; ngược lại diện tích lệch phải mạnh (độ lệch \SL{eda/house/area_skew}), sau
    $\log(1 + x)$ còn \SL{eda/house/logarea_skew}.
  \item \textbf{Tương quan yếu và phi tuyến:} tương quan tuyến tính của giá với diện tích chỉ $r = \SL{eda/house/corr_area}$,
    với số tầng $r = \SL{eda/house/corr_floors}$, với đường vào $r = \SL{eda/house/corr_road}$ (Hình~\ref{fig:house-corr}).
    Giá phụ thuộc mạnh vào vị trí: trung vị giá ở TP.~Hồ Chí Minh và Hà Nội cao nhất, Đà Nẵng thấp hơn một chút, còn Bình
    Dương và Đồng Nai thấp hơn gần một nửa (Hình~\ref{fig:house-scatter}), mà vị trí là biến phân loại có hàng trăm giá trị.
  \item \textbf{Thiếu dữ liệu nhiều:} \SL{eda/house/miss_frontagePct}\,\% tin thiếu mặt tiền, \SL{eda/house/miss_roadPct}\,\%
    thiếu đường vào, \SL{eda/house/miss_directionPct}\,\% thiếu hướng nhà.
  \item \textbf{Tiền xử lý:} tách tỉnh và quận/huyện từ địa chỉ; biến số được điền trung vị \emph{kèm cờ đánh dấu thiếu} rồi
    chuẩn hóa z-score; biến phân loại điền “Không rõ” rồi mã hóa one-hot, gộp các giá trị xuất hiện dưới 40 lần thành một
    nhóm “hiếm” để không sinh hàng trăm cột thưa.
\end{itemize}

\subsection{Tập dữ liệu 3: Phân loại cảm xúc nhận xét thương mại điện tử (\SL{eda/reviews/n} mẫu)}

\begin{itemize}
  \item \textbf{Nguồn:} Women's E-Commerce Clothing Reviews \cite{kaggle_reviews}, nhận xét thật (tiếng Anh) của khách
    hàng một nhà bán lẻ thời trang trực tuyến.
  \item \textbf{Quy mô:} \SL{eda/reviews/n_raw} bản ghi, giữ \SL{eda/reviews/n} bản có nội dung nhận xét. Nhãn cảm xúc ba
    lớp suy từ số sao: 1 -- 2 sao là \emph{tiêu cực}, 3 sao là \emph{trung tính}, 4 -- 5 sao là \emph{tích cực}. Tiêu đề và
    nội dung được ghép thành một văn bản.
\end{itemize}

\bang{tab_sample_rev}{Năm nhận xét tiêu biểu kèm nhãn cảm xúc.}{sample-rev}

\hinhpdf{f2_rev_dist}{0.95}{Phân bố số sao và nhãn cảm xúc: lớp tích cực chiếm \SL{eda/reviews/share_posPct}\,\%.}{rev-dist}
\hinhpdf{f2_rev_len}{0.65}{Phân bố độ dài nhận xét (số từ) theo nhãn; nét đứt là trung bình \SL{eda/reviews/words_mean} từ.}{rev-len}
\hinhpdf{f2_rev_words}{1}{Mười lăm từ đặc trưng nhất của nhận xét tích cực và tiêu cực, xếp theo tỉ số log-odds so với
phần còn lại.}{rev-words}

\paragraph{Nhận xét EDA và quy trình tiền xử lý văn bản.} Dữ liệu lệch mạnh về phía tích cực
(\SL{eda/reviews/share_posPct}\,\%), trong khi trung tính chỉ \SL{eda/reviews/share_neuPct}\,\% và tiêu cực
\SL{eda/reviews/share_negPct}\,\% (Hình~\ref{fig:rev-dist}); vì vậy chỉ số chính là F1 trung bình macro trên ba lớp.
Độ dài gần như không phân biệt được các lớp: trung bình \SL{eda/reviews/words_by_class/0} từ ở nhận xét tiêu cực,
\SL{eda/reviews/words_by_class/1} ở trung tính và \SL{eda/reviews/words_by_class/2} ở tích cực (Hình~\ref{fig:rev-len}),
nên tín hiệu nằm ở \emph{từ ngữ} chứ không ở độ dài. Từ đặc trưng nhất của lớp tiêu cực là \emph{\SL{eda/reviews/top_neg/0}},
\emph{\SL{eda/reviews/top_neg/1}}, \emph{\SL{eda/reviews/top_neg/2}}, \emph{\SL{eda/reviews/top_neg/3}}; của lớp tích cực là
\emph{\SL{eda/reviews/top_pos/3}}, \emph{\SL{eda/reviews/top_pos/4}}, \emph{\SL{eda/reviews/top_pos/0}}
(Hình~\ref{fig:rev-words}). Quy trình tiền xử lý:
\begin{enumerate}
  \item chuyển về chữ thường, bỏ URL và mọi ký tự không phải chữ cái hay dấu nháy;
  \item loại từ dừng tiếng Anh (danh sách của scikit-learn);
  \item vector hóa TF-IDF với unigram và bigram, giữ 5.000 đặc trưng tần suất cao nhất, bỏ từ xuất hiện dưới 3 văn bản,
    dùng $1 + \log(\text{tf})$ để giảm ảnh hưởng của từ lặp nhiều lần.
\end{enumerate}

% ---------------------------------------------------------------------
\section{Hiện thực hóa và thực nghiệm các mô hình học máy cơ bản}

Mỗi mô hình được đóng gói thành một \texttt{Pipeline} của scikit-learn \cite{pedregosa2011} gồm bộ tiền xử lý và bộ ước
lượng. Gọi \texttt{fit} trên phần huấn luyện sẽ khớp cả hai, nên không thống kê nào của tập kiểm thử rò vào quá trình
học; lưu pipeline bằng \texttt{joblib} là có ngay đối tượng nhận dữ liệu thô cho Web App.

\subsection{Pipeline hồi quy tuyến tính và Ridge}

Bộ tiền xử lý giá nhà (Đoạn mã~\ref{code:data_tabular__house_preprocessor}) xử lý riêng nhóm cột số và nhóm cột phân loại
bằng \texttt{ColumnTransformer}; Đoạn mã~\ref{code:ml_classic__build_linear_regression_pipeline} ghép nó với mô hình
tuyến tính, \texttt{alpha=None} cho bình phương tối thiểu thường và \texttt{alpha=1.0} cho Ridge.

\maNguon{data_tabular__house_preprocessor}{Bộ tiền xử lý dữ liệu giá nhà: điền thiếu kèm cờ, chuẩn hóa, one-hot có gộp nhóm hiếm}
\maNguon{ml_classic__build_linear_regression_pipeline}{Pipeline hồi quy tuyến tính và Ridge với scikit-learn}

\subsection{Pipeline hồi quy Logistic cho dữ liệu y tế}

\texttt{class\_weight="balanced"} bù cho tỉ lệ dương \SL{eda/diabetes/pos_ratePct}\,\%; bộ giải \texttt{lbfgs} là phương
pháp quasi-Newton hội tụ nhanh với số đặc trưng nhỏ.

\maNguon{ml_classic__build_logistic_regression_pipeline}{Pipeline hồi quy Logistic có cân bằng trọng số lớp}

\subsection{Pipeline Rừng ngẫu nhiên}

Rừng gồm 200 cây không giới hạn độ sâu, mỗi lá tối thiểu 2 mẫu; với phân loại, trọng số lớp được cân bằng lại trên từng
mẫu bootstrap. Cây không cần chuẩn hóa đặc trưng, nhưng dùng chung bộ tiền xử lý giúp mọi mô hình nhận cùng một đầu vào.

\maNguon{ml_classic__build_rf_pipeline}{Pipeline Random Forest cho cả hồi quy và phân loại}

\subsection{Pipeline Máy vector hỗ trợ}

Nhân RBF tính khoảng cách Euclid nên đặc trưng bắt buộc cùng thang đo; bộ tiền xử lý đã chuẩn hóa z-score. Văn bản dùng
\texttt{LinearSVC}.

\maNguon{ml_classic__build_svm_pipeline}{Pipeline SVM: SVC/SVR nhân RBF cho dữ liệu bảng, LinearSVC cho văn bản}

\subsection{Kết quả thực nghiệm các mô hình học máy cơ bản}

Ba bảng dưới đây tách theo bản chất bài toán. Thời gian là thời gian huấn luyện trên CPU (Random Forest dùng 12 luồng);
dung lượng là kích thước tệp \texttt{joblib} nén.

\bang{tab_ch2_diab_sk}{Kết quả các mô hình học máy cơ bản trên tập tiểu đường (\SL{eda/diabetes/n} hồ sơ, kiểm thử 20\,\%).}{ch2-diab-sk}
\bang{tab_ch2_house_sk}{Kết quả các mô hình học máy cơ bản trên tập giá nhà Việt Nam.}{ch2-house-sk}
\bang{tab_ch2_rev_sk}{Kết quả các mô hình học máy cơ bản trên tập nhận xét khách hàng (TF-IDF 5.000 chiều).}{ch2-rev-sk}

\subsection{Phân tích kết quả và giới hạn của các mô hình cơ bản}

\begin{enumerate}
  \item \textbf{Mô hình tuyến tính.} Trên tập tiểu đường, Logistic Regression có trọng số cân bằng đạt Recall
    \SL{ch2_sklearn/diabetes/logreg/recall} nhưng Precision chỉ \SL{ch2_sklearn/diabetes/logreg/precision}: trọng số lớp
    đẩy ngưỡng quyết định về phía bắt nhiều ca bệnh, đổi lại hơn một nửa cảnh báo là báo động giả. ROC-AUC
    \SL{ch2_sklearn/diabetes/logreg/auc} cho thấy khả năng \emph{xếp hạng} vẫn tốt; điểm yếu nằm ở ranh giới tuyến tính
    không bắt được ngưỡng gãy của HbA1c. Trên tập giá nhà, OLS và Ridge cùng chỉ đạt $R^2 \approx \SL{ch2_sklearn/house/ols/r2}$
    (MAE \SL{ch2_sklearn/house/ols/mae} tỷ): phạt L2 gần như không giúp gì vì vấn đề là thiếu khớp chứ không phải đa
    cộng tuyến.
  \item \textbf{SVM nhân RBF.} SVC nâng ROC-AUC lên \SL{ch2_sklearn/diabetes/svm/auc} và F1 lên
    \SL{ch2_sklearn/diabetes/svm/f1}, nhưng huấn luyện mất \SL{ch2_sklearn/diabetes/svm/train_s}\,s trên
    80.000 mẫu, gấp hàng trăm lần Logistic Regression, và suy luận chậm nhất (\SL{ch2_sklearn/diabetes/svm/infer_ms}\,ms
    mỗi mẫu) vì phải so mẫu mới với mọi vector hỗ trợ. SVR đạt $R^2 = \SL{ch2_sklearn/house/svr/r2}$ trên giá nhà. Với văn
    bản, LinearSVC huấn luyện trong \SL{ch2_sklearn/reviews/linsvm/train_s}\,s và đạt Accuracy
    \SL{ch2_sklearn/reviews/linsvm/accPct}\,\%.
  \item \textbf{Random Forest.} Là mô hình cơ bản mạnh nhất trên cả hai tập bảng: F1 \SL{ch2_sklearn/diabetes/rf/f1} trên
    tiểu đường và $R^2 = \SL{ch2_sklearn/house/rf/r2}$, MAE \SL{ch2_sklearn/house/rf/mae} tỷ trên giá nhà. Đổi lại, tệp mô
    hình nặng \SL{ch2_sklearn/house/rf/size_kb}\,KB (200 cây sâu), không ngoại suy được ra ngoài khoảng giá đã thấy, và trên
    văn bản 5.000 chiều thưa thì kém hơn mô hình tuyến tính (F1 macro \SL{ch2_sklearn/reviews/rf/f1}) vì mỗi lần chia
    chỉ xét một từ.
  \item \textbf{Động lực chuyển sang MLP.} Mô hình tuyến tính thiếu khả năng biểu diễn, SVM nhân không mở rộng theo số
    mẫu, Random Forest cồng kềnh. Mạng nơ-ron đa tầng hứa hẹn cả ba: xấp xỉ được hàm phi tuyến bất kỳ (định lý xấp xỉ phổ
    quát), huấn luyện theo lô nhỏ nên chi phí tuyến tính theo $N$ và chạy song song trên GPU, tệp mô hình chỉ vài chục
    đến vài trăm KB.
\end{enumerate}

% ---------------------------------------------------------------------
\section{Mạng nơ-ron đa tầng (MLP) trên ba nền tảng: NumPy, Keras, PyTorch}

\subsection{Cơ sở lý thuyết và kiến trúc mạng}

Một MLP gồm các tầng kết nối đầy đủ. Tầng $l$ nhận kích hoạt $\mathbf{a}^{[l-1]}$ của tầng trước và tính
\begin{equation}
  \mathbf{z}^{[l]} = W^{[l]}\mathbf{a}^{[l-1]} + \mathbf{b}^{[l]}, \qquad \mathbf{a}^{[l]} = g^{[l]}\big(\mathbf{z}^{[l]}\big),
\end{equation}
với $g$ là ReLU ở tầng ẩn và sigmoid, softmax hoặc hàm đồng nhất ở tầng ra tùy bài toán. Lan truyền ngược áp dụng quy tắc
dây chuyền từ tầng ra về tầng vào:
\begin{align}
  \boldsymbol{\delta}^{[L]} &= \hat{\mathbf{y}} - \mathbf{y}, \\
  \boldsymbol{\delta}^{[l]} &= \big((W^{[l+1]})^\top\boldsymbol{\delta}^{[l+1]}\big) \odot g'^{[l]}\big(\mathbf{z}^{[l]}\big), \\
  \frac{\partial\mathcal{L}}{\partial W^{[l]}} &= \boldsymbol{\delta}^{[l]}\big(\mathbf{a}^{[l-1]}\big)^\top, \qquad
  \frac{\partial\mathcal{L}}{\partial\mathbf{b}^{[l]}} = \boldsymbol{\delta}^{[l]} .
\end{align}
Đẳng thức đầu đúng cho cả ba cặp đầu ra -- mất mát dùng ở đây (sigmoid + BCE, softmax + entropy chéo, tuyến tính + MSE
với hệ số 2), nên một hàm \texttt{backward} phục vụ cả ba bài toán. Trọng số khởi tạo theo He \cite{he2015init}
($\mathcal{N}(0, 2/n_{in})$) để biên độ kích hoạt ổn định qua các tầng ReLU, và cập nhật bằng Adam \cite{kingma2015}.

Để so sánh ba nền tảng một cách công bằng, cả ba cài đặt dùng \emph{cùng một kiến trúc} (Hình~\ref{fig:mlp}): hai tầng ẩn
128 và 64 nơ-ron, ReLU, Dropout 0,2 \cite{srivastava2014}; Adam với tốc độ học $10^{-3}$, lô 256; dừng sớm khi mất mát
kiểm định không giảm trong 8 epoch và khôi phục trọng số tốt nhất.

\hinhtikz{f2_mlp}{Kiến trúc MLP dùng chung cho cả ba nền tảng và ba bài toán.}{mlp}

\subsection{Hiện thực từ đầu bằng NumPy}

Lớp \texttt{MLPScratch} chỉ dùng phép toán ma trận của NumPy. Lan truyền tiến (Đoạn mã~\ref{code:mlp_scratch__MLPScratch.forward})
lưu lại $\mathbf{z}$, $\mathbf{a}$ và mặt nạ Dropout của từng tầng; lan truyền ngược
(Đoạn mã~\ref{code:mlp_scratch__MLPScratch.backward}) dùng lại chúng để tính gradient đúng như các công thức trên; Adam
(Đoạn mã~\ref{code:mlp_scratch__MLPScratch._adam}) giữ hai mômen cho mỗi tham số và hiệu chỉnh độ chệch khởi đầu.

\maNguon{mlp_scratch__MLPScratch.forward}{Lan truyền tiến của MLP viết tay bằng NumPy, kèm Dropout ngược}
\maNguon{mlp_scratch__MLPScratch.backward}{Lan truyền ngược của MLP viết tay: một sai số tầng ra chung cho ba bài toán}
\maNguon{mlp_scratch__MLPScratch._adam}{Bộ tối ưu Adam viết tay}

\subsection{Hiện thực bằng Keras / TensorFlow}

Keras gói cả kiến trúc, hàm mất mát và vòng huấn luyện vào API cấp cao; phần đầu ra chọn theo bảng \texttt{HEADS}.

\maNguon{mlp_keras__build_keras_mlp}{MLP bằng Keras Sequential API cho ba loại đầu ra}

\subsection{Hiện thực bằng PyTorch}

PyTorch tách mô hình (\texttt{nn.Module}) khỏi vòng huấn luyện do ta tự viết. Mô hình trả logit thô; sigmoid và softmax
nằm trong \texttt{BCEWithLogitsLoss} và \texttt{CrossEntropyLoss} để ổn định số học. Vì dữ liệu bảng nhỏ, toàn bộ tập
được đưa lên GPU một lần và chia lô bằng chỉ số thay vì \texttt{DataLoader}.

\maNguon{mlp_torch__TorchMLP}{MLP bằng PyTorch với lớp \texttt{nn.Module}}
\maNguon{mlp_torch__train_epoch}{Một epoch huấn luyện PyTorch trên dữ liệu đã nằm sẵn trên GPU}

% ---------------------------------------------------------------------
\section{Đánh giá và đối chuẩn trên ba tập dữ liệu}

\subsection{Đối chuẩn hiệu năng chi tiết trên từng tập dữ liệu}

Ba bảng sau đặt ba bản MLP cạnh các mô hình cơ bản, trên cùng tập kiểm thử.

\bang{tab_ch2_diab_all}{Đối chuẩn mọi mô hình trên tập tiểu đường.}{ch2-diab-all}
\bang{tab_ch2_house_all}{Đối chuẩn mọi mô hình trên tập giá nhà Việt Nam.}{ch2-house-all}
\bang{tab_ch2_rev_all}{Đối chuẩn mọi mô hình trên tập nhận xét khách hàng.}{ch2-rev-all}
\bang{tab_ch2_fw}{Cùng một kiến trúc MLP trên ba nền tảng.}{ch2-fw}

\hinhpdf{f2_loss_mlp}{1}{Đường cong mất mát huấn luyện (nét đứt) và kiểm định (nét liền) của MLP trên ba nền tảng.}{loss-mlp}
\hinhpdf{f2_house_pred}{0.95}{Giá thực tế và giá dự báo trên tập kiểm thử của Random Forest và MLP PyTorch.}{house-pred}
\hinhpdf{f2_bar_diab}{1}{So sánh đa tiêu chí trên tập tiểu đường.}{bar-diab}
\hinhpdf{f2_bar_house}{1}{So sánh sai số (MAE, RMSE) và hệ số xác định $R^2$ trên tập giá nhà.}{bar-house}
\hinhpdf{f2_bar_rev}{1}{So sánh Accuracy và các chỉ số macro trên tập nhận xét khách hàng.}{bar-rev}
\hinhpdf{f2_cm_diab}{1}{Ma trận nhầm lẫn trên tập kiểm thử tiểu đường: Logistic Regression, Random Forest, MLP PyTorch.}{cm-diab}
\hinhpdf{f2_cm_rev}{1}{Ma trận nhầm lẫn trên tập kiểm thử nhận xét khách hàng.}{cm-rev}

\paragraph{Thảo luận kết quả.}
\begin{itemize}
  \item \textbf{MLP tốt nhất trên dữ liệu y tế.} Cả ba bản MLP đạt F1 lớp dương quanh
    \SL{ch2_mlp_torch/diabetes/f1} (PyTorch), cao hơn Random Forest (\SL{ch2_sklearn/diabetes/rf/f1}) và vượt xa Logistic
    Regression (\SL{ch2_sklearn/diabetes/logreg/f1}). Ma trận nhầm lẫn (Hình~\ref{fig:cm-diab}) cho thấy MLP gần như không
    báo động giả (Precision \SL{ch2_mlp_torch/diabetes/precision}) nhưng vẫn bỏ sót khoảng một phần ba ca bệnh
    (Recall \SL{ch2_mlp_torch/diabetes/recall}). ROC-AUC của bốn mô hình phi tuyến chỉ chênh nhau ở chữ số thứ ba
    (\SL{ch2_sklearn/diabetes/svm/auc} -- \SL{ch2_mlp_torch/diabetes/auc}): khác biệt lớn về F1 phần nhiều đến từ
    \emph{điểm cắt} trên đường ROC chứ không từ khả năng phân biệt. Trong sàng lọc y tế, bỏ sót đắt hơn báo động giả, nên
    khi triển khai nên hạ ngưỡng của MLP xuống dưới 0,5 để đổi Precision lấy Recall.
  \item \textbf{Random Forest vẫn nhỉnh nhất trên giá nhà.} $R^2$ của Random Forest (\SL{ch2_sklearn/house/rf/r2}) cao hơn
    MLP PyTorch (\SL{ch2_mlp_torch/house/r2}) một chút, hai mô hình bám nhau trên đồ thị phân tán
    (Hình~\ref{fig:house-pred}). Cả hai đều chỉ giải thích khoảng 62\,\% phương sai: phần còn lại nằm ở thông tin dữ liệu
    không có (vị trí chính xác trong quận, dự án, chất lượng xây dựng), nên không mô hình nào vượt xa được. Cây quyết
    định chia trực tiếp không gian theo từng quận/huyện, một lợi thế tự nhiên trên dữ liệu bảng có nhiều biến phân loại.
  \item \textbf{Văn bản: MLP chính xác nhất nhưng lớp trung tính là điểm mù chung.} MLP đạt Accuracy cao nhất
    (\SL{ch2_mlp_scratch/reviews/accPct}\,\% với bản NumPy), F1 macro ngang các mô hình tuyến tính. Cả ba bản MLP quá khớp
    rất nhanh: mất mát kiểm định chạm đáy ở epoch 2 -- 3 rồi tăng (Hình~\ref{fig:loss-mlp}), vì \SL{ch2_mlp_torch/reviews/params} tham số học
    trên chưa đến 17.000 văn bản. Ma trận nhầm lẫn (Hình~\ref{fig:cm-rev}) cho thấy lớp trung tính (3 sao) bị lẫn
    sang cả hai phía: một nhận xét 3 sao thường vừa khen vừa chê.
  \item \textbf{Ba nền tảng cho kết quả tương đương.} Với cùng kiến trúc và siêu tham số, chênh lệch F1/$R^2$ giữa NumPy,
    Keras và PyTorch nằm trong khoảng dao động do khởi tạo ngẫu nhiên khác nhau (Bảng~\ref{tab:ch2-fw}). Bản NumPy tự viết
    đạt độ chính xác ngang thư viện công nghiệp, xác nhận các công thức lan truyền ngược được cài đặt đúng. Khác biệt nằm
    ở chi phí kỹ thuật: PyTorch chạy GPU, Keras có tệp lưu lớn nhất do chứa cả trạng thái bộ tối ưu.
\end{itemize}

% ---------------------------------------------------------------------
\section{Đóng gói triển khai và giao diện Web thực tế}
\label{sec:web}

\subsection{Kiến trúc triển khai hệ thống}

Toàn bộ mô hình của ba chương được phục vụ bởi một ứng dụng Flask duy nhất gồm ba tầng (Hình~\ref{fig:deploy}):
\begin{enumerate}
  \item \textbf{Tầng giao diện:} một trang HTML5/CSS3 với sáu thẻ, JavaScript gọi API bằng \texttt{fetch} và vẽ biểu đồ
    bằng SVG thuần, không phụ thuộc thư viện ngoài.
  \item \textbf{Tầng ứng dụng:} máy chủ Flask nhận yêu cầu \texttt{POST} JSON tại sáu endpoint
    \texttt{/api/diabetes}, \texttt{/api/house}, \texttt{/api/sentiment}, \texttt{/api/vision}, \texttt{/api/stock},
    \texttt{/api/churn}.
  \item \textbf{Tầng mô hình:} mọi mô hình (\texttt{.pt}, \texttt{.joblib}) được nạp lên RAM một lần lúc khởi động và suy
    luận trên CPU. Bộ tiền xử lý được dựng lại bằng chính hàm \texttt{load\_*} đã dùng khi huấn luyện nên dữ liệu người
    dùng đi qua đúng phép biến đổi mà mô hình đã học.
\end{enumerate}

\hinhtikz{f2_deploy}{Kiến trúc ba tầng của Web App phục vụ sáu phân hệ thông minh.}{deploy}

Endpoint sàng lọc tiểu đường (Đoạn mã~\ref{code:app__api_diabetes}) minh họa khuôn chung: dựng một dòng
\texttt{DataFrame} từ JSON, biến đổi bằng bộ tiền xử lý, chạy mô hình, đo thời gian và trả JSON.

\maNguon{app__api_diabetes}{Endpoint sàng lọc tiểu đường: MLP PyTorch và Random Forest cho cùng một hồ sơ}

\subsection{Giao diện thực tế các hệ thống thông minh}

\anhUI{ui_diabetes}{Phân hệ sàng lọc nguy cơ tiểu đường: người dùng nhập chỉ số lâm sàng, hệ thống trả xác suất của MLP và
Random Forest, mức nguy cơ và khuyến nghị.}{ui-diabetes}
\anhUI{ui_house}{Phân hệ định giá bất động sản: giá ước tính của Random Forest, khoảng hợp lý $\pm$ RMSE kiểm thử, dự báo
của MLP và đơn giá theo m².}{ui-house}
\anhUI{ui_sentiment}{Phân hệ phân tích cảm xúc nhận xét: xác suất ba lớp của MLP, đối chiếu Logistic Regression và các từ
đóng góp nhiều nhất.}{ui-sentiment}

Ba phân hệ trả kết quả trong vài đến vài chục mili-giây trên CPU, cho thấy các mô hình học máy cơ bản và MLP chuyển được trực tiếp
vào một quy trình nghiệp vụ thực tế.

% =====================================================================
%  CHƯƠNG 3 — MẠNG NƠ-RON TÍCH CHẬP (CNN)
% =====================================================================
\chapter{Mạng Nơ-ron Tích Chập (CNN)}
\label{ch:cnn}

\section{Cơ sở lý thuyết mạng nơ-ron tích chập và biểu diễn không gian}

MLP của Chương~\ref{ch:hoc-may} gặp hai trở ngại khi nhận đầu vào là ảnh:
\begin{enumerate}
  \item \textbf{Bùng nổ tham số.} Duỗi một ảnh màu $224{\times}224{\times}3$ thành vector 150.528 chiều, chỉ một tầng
    ẩn 1.024 nơ-ron đã cần hơn 154 triệu trọng số, vừa tốn bộ nhớ vừa dễ quá khớp.
  \item \textbf{Mất cấu trúc không gian.} Duỗi phẳng xóa thông tin “điểm nào cạnh điểm nào”; cùng một nét chữ dịch sang
    phải hai điểm ảnh trở thành một vector hoàn toàn khác.
\end{enumerate}
Mạng nơ-ron tích chập \cite{lecun1998} khắc phục bằng ba nguyên lý: \emph{vùng tiếp nhận cục bộ} (mỗi nơ-ron chỉ nhìn một
ô nhỏ), \emph{chia sẻ trọng số} (cùng một bộ lọc quét toàn ảnh, nên một mẫu hình được nhận ra ở mọi vị trí) và
\emph{lấy mẫu giảm} (pooling). Đầu vào được biểu diễn thành tensor: ảnh xám là ma trận $H{\times}W$, ảnh màu là ba ma trận
R, G, B xếp chồng thành khối $C{\times}H{\times}W$, và một lô $N$ ảnh là tensor bốn chiều (Hình~\ref{fig:tensor3}).

\hinhtikz{f3_tensor}{Biểu diễn ảnh số: từ ma trận điểm ảnh đến tensor ba kênh và lô ảnh 4D trong học sâu.}{tensor3}

\subsection{Toán học của phép tích chập hai chiều}

Trong học sâu, “tích chập” thực chất là tương quan chéo. Với tensor vào $X \in \mathbb{R}^{C_{in}\times H\times W}$ và
$C_{out}$ bộ lọc $K \in \mathbb{R}^{C_{out}\times C_{in}\times k_h\times k_w}$, giá trị tại vị trí $(i, j)$ của kênh ra $c$ là
\begin{equation}
  Y_{c,i,j} = \sum_{m=1}^{C_{in}}\sum_{u=0}^{k_h-1}\sum_{v=0}^{k_w-1} X_{m,\,i s+u,\,j s+v}\,K_{c,m,u,v} + b_c ,
  \label{eq:conv}
\end{equation}
với bước trượt $s$. Mỗi bộ lọc trải hết chiều sâu $C_{in}$ của khối vào và trượt theo hai chiều không gian; mỗi vị trí cho
\emph{một} số, nên một bộ lọc sinh ra một kênh ra, và $C_{out}$ bộ lọc sinh ra khối ra $C_{out}$ kênh
(Hình~\ref{fig:conv3d}). Kích thước bản đồ đặc trưng ra phụ thuộc đệm số 0 $p$:
\begin{equation}
  H_{out} = \Big\lfloor\frac{H - k_h + 2p}{s}\Big\rfloor + 1, \qquad W_{out} = \Big\lfloor\frac{W - k_w + 2p}{s}\Big\rfloor + 1 .
  \label{eq:outsize}
\end{equation}

\hinhtikz{f3_conv3d}{Phép tích chập nhìn trong không gian 3D: bộ lọc $C_{in}{\times}k{\times}k$ trượt trên khối vào, mỗi vị
trí cho một phần tử của khối ra.}{conv3d}

Hình~\ref{fig:conv} tính tay một ví dụ với bộ lọc phát hiện cạnh dọc. Số tham số của một tầng chỉ là
$C_{out}(C_{in}k_hk_w + 1)$, không phụ thuộc kích thước ảnh: tầng Conv $3{\times}3$ từ 64 lên 128 kênh có
$128 \cdot (64 \cdot 9 + 1) = 73.856$ tham số dù ảnh vào lớn đến đâu.

\hinhtikz{f3_conv}{Minh họa số học của phép tích chập 2D: ảnh $5{\times}5$, bộ lọc cạnh dọc $3{\times}3$, bước 1; ô đậm
màu ứng với giá trị lớn.}{conv}

\subsection{Vùng cảm nhận và phép gộp (pooling)}

Vùng cảm nhận (receptive field, RF) là vùng ảnh vào mà một nơ-ron ở tầng sâu “nhìn thấy”. Xếp chồng các tầng $k{\times}k$
làm RF tăng dần:
\begin{equation}
  RF_l = RF_{l-1} + (k_l - 1)\prod_{i=1}^{l-1} s_i .
\end{equation}
Với bước 1, ba tầng $3{\times}3$ cho RF $= 7$ (Hình~\ref{fig:rf}); mỗi lần pooling bước 2 nhân đôi tốc độ tăng của RF ở
các tầng sau. Nhờ vậy các tầng đầu học nét và cạnh, các tầng giữa học kết cấu và chi tiết, các tầng cuối nhìn được toàn bộ
đối tượng.

\hinhtikz{f3_rf}{Cơ chế tích lũy vùng cảm nhận khi xếp chồng ba tầng tích chập $3{\times}3$ bước 1.}{rf}

Pooling giảm kích thước không gian, tăng khả năng chịu dịch chuyển nhỏ và mở rộng nhanh vùng cảm nhận
(Hình~\ref{fig:pool}). \textbf{Max pooling} giữ giá trị lớn nhất trong cửa sổ (đặc trưng mạnh nhất);
\textbf{average pooling} lấy trung bình. \textbf{Global Average Pooling} (GAP) \cite{lin2013} lấy trung bình toàn bộ mỗi
bản đồ đặc trưng, biến $C{\times}H{\times}W$ thành vector $C$ chiều, thay cho tầng Flatten + Dense khổng lồ. Mọi kiến trúc
trong chương này kết thúc bằng GAP, nên cùng một mạng nhận được cả ảnh $32{\times}32$ lẫn $64{\times}64$.

\hinhtikz{f3_pool}{So sánh max pooling, average pooling ($2{\times}2$, bước 2) và Global Average Pooling trên cùng một
bản đồ đặc trưng.}{pool}

\subsection{Tiến hóa kiến trúc: VGG, ResNet và MobileNet}

Hình~\ref{fig:evolution} tóm tắt ba thế hệ khối xây dựng được hiện thực trong chương này.

\hinhtikz{f3_evolution}{Sơ đồ tiến hóa khối xây dựng CNN: khối VGG phẳng, khối Residual và khối Inverted Residual.}{evolution}

\paragraph{VGG \cite{simonyan2015}.} Các khối gồm hai hoặc ba tầng Conv $3{\times}3$ rồi MaxPool, số kênh nhân đôi sau mỗi
lần giảm kích thước. Hai tầng $3{\times}3$ nhìn được vùng $5{\times}5$ với 18 thay vì 25 trọng số mỗi cặp kênh và có hai
hàm ReLU thay vì một (Hình~\ref{fig:vgg}). Thiết kế đơn giản, đều đặn nhưng nhiều tham số và khó tối ưu khi rất sâu.

\hinhtikz{f3_vgg}{Khối VGG: hai tầng $3{\times}3$ có cùng vùng nhìn $5{\times}5$ với một tầng $5{\times}5$ nhưng ít trọng
số hơn.}{vgg}

\paragraph{ResNet \cite{he2016}.} Khi mạng sâu quá khoảng 20 tầng, sai số \emph{huấn luyện} lại tăng (thoái hóa), không
phải do quá khớp mà do tối ưu khó. Khối Residual học phần dư $F(\mathbf{x})$ rồi cộng với đầu vào,
$\mathbf{y} = F(\mathbf{x}) + \mathbf{x}$. Theo quy tắc dây chuyền
\begin{equation}
  \frac{\partial E}{\partial \mathbf{x}} = \frac{\partial E}{\partial \mathbf{y}}\Big(\frac{\partial F}{\partial \mathbf{x}} + I\Big),
\end{equation}
ma trận đơn vị $I$ là “đường cao tốc” để gradient về tới tầng đầu ngay cả khi $\partial F/\partial\mathbf{x}$ nhỏ
(Hình~\ref{fig:residual}).

\hinhtikz{f3_residual}{Khối Residual: lan truyền xuôi $y = F(x) + x$ và dòng gradient qua kết nối tắt.}{residual}

\paragraph{MobileNet \cite{howard2017,sandler2018}.} Tích chập tách biệt chiều sâu chia tích chập chuẩn thành
\emph{depthwise} (mỗi kênh một bộ lọc $k{\times}k$, chi phí $C_{in}HWk^2$) và \emph{pointwise} (Conv $1{\times}1$ trộn
kênh, chi phí $C_{in}C_{out}HW$). Tỉ lệ chi phí so với tích chập chuẩn là
\begin{equation}
  \frac{C_{in}HWk^2 + C_{in}C_{out}HW}{C_{in}C_{out}HWk^2} = \frac{1}{C_{out}} + \frac{1}{k^2} \approx \frac{1}{9}\quad (k = 3).
\end{equation}
MobileNetV2 bọc phép này trong khối \emph{Inverted Residual}: Conv $1{\times}1$ mở rộng số kênh lên 6 lần, depthwise
$3{\times}3$ ở không gian rộng, rồi Conv $1{\times}1$ chiếu về ít kênh \emph{không} qua ReLU (nút cổ chai tuyến tính, tránh
mất thông tin ở không gian hẹp), kèm kết nối tắt khi kích thước giữ nguyên (Hình~\ref{fig:depthwise}).

\hinhtikz{f3_depthwise}{Tích chập chuẩn so với tích chập tách biệt chiều sâu: lọc không gian riêng từng kênh rồi trộn
kênh bằng $1{\times}1$.}{depthwise}

% ---------------------------------------------------------------------
\section{Khảo sát thực nghiệm và tiền xử lý ba bộ dữ liệu đa miền}

\subsection{Tập dữ liệu 1: nhận dạng chữ số nhà SVHN}

\begin{itemize}
  \item \textbf{Nguồn:} Street View House Numbers của Stanford \cite{netzer2011}, cắt từ ảnh Google Street View; tải qua
    lớp \texttt{SVHN} của \texttt{torchvision}.
  \item \textbf{Quy mô:} \SL{eda/svhn/n_train} ảnh huấn luyện và \SL{eda/svhn/n_test} ảnh kiểm thử theo cách chia gốc, ảnh
    màu $32{\times}32{\times}3$, 10 lớp chữ số 0 -- 9.
\end{itemize}

\hinhpdf{f3_svhn_grid}{1}{Bốn ảnh mẫu mỗi lớp của SVHN: đa dạng phông chữ, ánh sáng, màu nền và có chữ số lân cận lọt vào
khung.}{svhn-grid}
\hinhpdf{f3_svhn_classes}{1}{Số ảnh theo lớp của SVHN trong tập huấn luyện và kiểm thử.}{svhn-classes}
\hinhpdf{f3_svhn_rgb}{1}{Phân bố cường độ điểm ảnh trên ba kênh màu của SVHN.}{svhn-rgb}
\hinhpdf{f3_svhn_mean}{1}{Ảnh trung bình của từng lớp SVHN (đã kéo giãn tương phản): hình dáng chữ số vẫn hiện rõ dù nền
và màu rất đa dạng.}{svhn-mean}

\paragraph{Nhận xét EDA.} Khác MNIST viết tay trên nền trắng, ảnh SVHN có nền thật (tường, cửa, biển kim loại), độ tương
phản thấp và thường có chữ số hàng xóm lọt vào hai bên (Hình~\ref{fig:svhn-grid}). Phân bố lớp lệch: chữ số 1 nhiều nhất
(\SL{eda/svhn/max_class_n} ảnh), chữ số 9 ít nhất (\SL{eda/svhn/min_class_n} ảnh), vì số nhà bắt đầu bằng 1 nhiều hơn
(Hình~\ref{fig:svhn-classes}). Ba kênh có phân bố gần nhau (Hình~\ref{fig:svhn-rgb}), trung bình
$(\SL{eda/svhn/mean/0}; \SL{eda/svhn/mean/1}; \SL{eda/svhn/mean/2})$ và độ lệch chuẩn
$(\SL{eda/svhn/std/0}; \SL{eda/svhn/std/1}; \SL{eda/svhn/std/2})$, dùng để chuẩn hóa z-score từng kênh: màu sắc gần như
không mang thông tin về nhãn, hình dáng mới là tín hiệu. Ảnh trung bình mỗi lớp (Hình~\ref{fig:svhn-mean}) cho thấy chữ số
luôn được căn giữa khung. Tăng cường dữ liệu chỉ dùng dịch ngẫu nhiên $\pm 4$ điểm ảnh (đệm phản chiếu); \emph{không} lật
ảnh vì chữ số lật không còn là chữ số đó.

\subsection{Tập dữ liệu 2: phân loại ảnh vệ tinh EuroSAT}

\begin{itemize}
  \item \textbf{Nguồn:} EuroSAT \cite{helber2019}, ảnh vệ tinh Sentinel-2 của Cơ quan Vũ trụ châu Âu, bản RGB, tải qua
    lớp \texttt{EuroSAT} của \texttt{torchvision}.
  \item \textbf{Quy mô:} 27.000 ảnh $64{\times}64{\times}3$, 10 lớp sử dụng đất (cây hằng năm, rừng, thảm cỏ, đường cao
    tốc, khu công nghiệp, đồng cỏ chăn thả, cây lâu năm, khu dân cư, sông, biển/hồ). Không có tập kiểm thử chính thức nên
    chia phân tầng 80/20: \SL{eda/eurosat/n_train} ảnh huấn luyện, \SL{eda/eurosat/n_test} ảnh kiểm thử.
\end{itemize}

\hinhpdf{f3_eurosat_grid}{1}{Bảy ảnh mẫu của mỗi lớp EuroSAT (mỗi hàng một lớp).}{eurosat-grid}
\hinhpdf{f3_eurosat_classes}{1}{Số ảnh theo lớp của EuroSAT.}{eurosat-classes}
\hinhpdf{f3_eurosat_rgb}{1}{Phân bố cường độ điểm ảnh trên ba kênh màu của EuroSAT.}{eurosat-rgb}
\hinhpdf{f3_eurosat_mean}{1}{Ảnh trung bình của từng lớp EuroSAT: biển/hồ, rừng và sông có màu trung bình đặc trưng.}{eurosat-mean}

\paragraph{Nhận xét EDA.} Các lớp gần cân bằng, từ \SL{eda/eurosat/min_class_n} đến \SL{eda/eurosat/max_class_n} ảnh
huấn luyện mỗi lớp (Hình~\ref{fig:eurosat-classes}). Phân biệt lớp dựa nhiều vào \emph{kết cấu} (vân ruộng, lưới đường phố)
và \emph{màu} (nước xanh thẫm, rừng xanh tối) hơn là hình dạng đối tượng: ảnh trung bình của biển/hồ gần như một mảng xanh
đồng nhất, trong khi khu dân cư và khu công nghiệp có ảnh trung bình xám mờ (Hình~\ref{fig:eurosat-mean}). Ảnh chụp từ trên
xuống không có chiều “trên dưới”, nên ngoài dịch ngẫu nhiên còn lật ngang, lật dọc và chuyển vị (tương đương xoay $90^\circ$)
mà không đổi nhãn.

\subsection{Tập dữ liệu 3: chuỗi đặc trưng lâm sàng cho 1D CNN}

Để thấy tích chập không chỉ dành cho ảnh, ta áp dụng 1D CNN lên tập tiểu đường của Chương~\ref{ch:hoc-may}. Mười lăm đặc
trưng sau tiền xử lý (4 biến số chuẩn hóa, 2 biến nhị phân, 9 cột one-hot) được xếp thành tensor $[B, 1, 15]$; hai tầng
Conv1D kích thước $1{\times}3$ (32 và 64 bộ lọc) học tương tác giữa các đặc trưng kề nhau, rồi tầng Dense 64 và đầu ra
sigmoid, tổng cộng \SL{ch3_diabetes1d/cnn1d/params} tham số. Mô hình dùng cùng cách chia và cùng dừng sớm như MLP.

% ---------------------------------------------------------------------
\section{Hiện thực hóa mã nguồn trên ba nền tảng}

\subsection{Tích chập 2D và Max Pooling viết tay bằng NumPy}

Đoạn mã~\ref{code:cnn_scratch__conv2d_loop} cài đặt đúng công thức~\eqref{eq:conv} bằng bốn vòng lặp để dễ đối chiếu.
Cách này quá chậm cho cả mạng, nên Đoạn mã~\ref{code:cnn_scratch__conv2d_im2col} dùng \texttt{sliding\_window\_view} để
lấy mọi cửa sổ thành một tensor rồi nhân một lần bằng \texttt{einsum}, cũng là ý tưởng im2col mà các thư viện thật dùng.

\maNguon{cnn_scratch__conv2d_loop}{Tích chập 2D nhiều kênh bằng vòng lặp, bám sát công thức (3.1)}
\maNguon{cnn_scratch__conv2d_im2col}{Tích chập 2D kiểu im2col bằng NumPy}
\maNguon{cnn_scratch__maxpool2d}{Max Pooling bằng NumPy}

Để kiểm chứng, ta nạp trọng số của Custom CNN đã huấn luyện bằng PyTorch trên SVHN vào các hàm NumPy và chạy toàn bộ mạng
(ba khối conv -- BN -- ReLU, GAP, Dense) trên \SL{ch3_scratch/n_img} ảnh kiểm thử:
\begin{itemize}
  \item bản vòng lặp và bản im2col lệch nhau tối đa \SL{ch3_scratch/loop_vs_im2col_maxdiff} (sai số làm tròn), nhưng
    tốn \SL{ch3_scratch/loop_ms_per_img}\,ms so với \SL{ch3_scratch/im2col_ms_per_img}\,ms mỗi ảnh cho riêng tầng đầu;
  \item logit của mạng NumPy lệch PyTorch tối đa \SL{ch3_scratch/net_maxdiff}; nhãn dự đoán trùng
    \SL{ch3_scratch/agreePct}\,\%, độ chính xác \SL{ch3_scratch/acc_numpyPct}\,\% so với \SL{ch3_scratch/acc_torchPct}\,\%;
  \item cả mạng NumPy mất \SL{ch3_scratch/numpy_net_ms_per_img}\,ms mỗi ảnh, PyTorch trên cùng CPU
    \SL{ch3_scratch/torch_cpu_ms_per_img}\,ms: hiểu thuật toán thì viết tay được, nhưng hiệu năng cần thư viện tối ưu.
\end{itemize}

\subsection{Custom CNN với khối kiểu VGG bằng Keras / TensorFlow}

Mô hình cơ sở gồm ba khối kiểu VGG (32, 64, 128 kênh), mỗi khối hai Conv $3{\times}3$ kèm BatchNorm \cite{ioffe2015} và
ReLU, MaxPool, Dropout tăng dần; sau đó GAP, Dense 128 và softmax (Hình~\ref{fig:arch-custom}). Cùng kiến trúc này được viết
lại bằng PyTorch để so sánh hai nền tảng.

\hinhtikz{f3_arch_custom}{Kiến trúc Custom CNN trên ảnh $32{\times}32$: bề rộng khối tỉ lệ với số kênh, chiều cao và chiều
sâu tỉ lệ với kích thước không gian.}{arch-custom}

\maNguon{cnn_keras__build_vgg_block}{Khối tích chập kiểu VGG bằng Keras}
\maNguon{cnn_keras__build_keras_custom_cnn}{Custom CNN bằng Keras Functional API}

\subsection{Khối Inverted Residual kiểu MobileNetV2 bằng PyTorch}

\maNguon{cnn_torch__InvertedResidual}{Khối Inverted Residual: mở rộng $1{\times}1$, depthwise $3{\times}3$, chiếu tuyến tính $1{\times}1$}

Toàn mạng MobileNetV2 xếp 17 khối theo bảng cấu hình của bài báo gốc, đổi bước của hai nhóm khối đầu từ 2 về 1 vì ảnh vào
nhỏ hơn ImageNet nhiều lần. Với ảnh EuroSAT $64{\times}64$, tầng thân giữ bước 2 như bản gốc để bản đồ đặc trưng sau thân là
$32{\times}32$, đúng bằng trường hợp SVHN. Lần chạy đầu giữ nguyên $64{\times}64$ qua các khối mở rộng 6 lần đã cần hơn
8\,GB bộ nhớ GPU ở lô 256: VRAM tràn sang bộ nhớ chia sẻ của Windows và một epoch kéo dài hơn 35 phút, nên cấu hình này được
thay.

\subsection{Kiến trúc ResNet-18 bằng PyTorch}

Hình~\ref{fig:arch-resnet} là toàn cảnh ResNet-18 cho ảnh $32{\times}32$: thân Conv $3{\times}3$ giữ nguyên độ phân giải,
bốn tầng mỗi tầng hai khối Residual, mỗi lần sang tầng mới kích thước không gian giảm một nửa và số kênh tăng gấp đôi.

\hinhtikz{f3_arch_resnet}{Kiến trúc ResNet-18 cho ảnh $32{\times}32$; mỗi cung cam là một kết nối tắt quanh một khối Residual.}{arch-resnet}

\maNguon{cnn_torch__ResidualBlock}{Khối Residual cơ bản với kết nối tắt chiếu $1{\times}1$ khi đổi kích thước}
\maNguon{cnn_torch__ResNet18}{Kiến trúc ResNet-18: thân $3{\times}3$, bốn tầng $[2, 2, 2, 2]$ khối, GAP và FC}

\subsection{Quy trình huấn luyện trên GPU}

Mọi mô hình PyTorch huấn luyện từ đầu với SGD Nesterov (momentum 0,9, weight decay $5\cdot10^{-4}$), lịch học One-Cycle
\cite{smith2019} đỉnh 0,1, lô 256; SVHN 15 epoch, EuroSAT 30 epoch. Toàn bộ ảnh \texttt{uint8} nạp lên GPU một lần; chuẩn
hóa và tăng cường dữ liệu làm ngay trên GPU (Đoạn mã~\ref{code:run_ch3_torch__augment}), và lan truyền chạy ở độ chính xác
hỗn hợp bfloat16. Nhờ vậy CPU rảnh cho các tiến trình Keras và scikit-learn chạy song song.

\maNguon{run_ch3_torch__augment}{Tăng cường dữ liệu trên GPU: dịch ngẫu nhiên từng ảnh bằng chỉ số nâng cao, lật và chuyển vị cho ảnh vệ tinh}

% ---------------------------------------------------------------------
\section{Đánh giá và đối chuẩn các biến thể kiến trúc CNN}

\paragraph{Phương pháp.} Bốn kiến trúc được huấn luyện từ trọng số ngẫu nhiên trên cùng dữ liệu, cùng siêu tham số, không
dùng trọng số tiền huấn luyện ImageNet, để phản ánh năng lực trích đặc trưng của từng cấu trúc trên dữ liệu đích:
\begin{enumerate}
  \item \textbf{Custom CNN} (mô hình cơ sở): 6 tầng Conv trong ba khối kiểu VGG và GAP, \SL{ch3_svhn/custom/params} tham số.
  \item \textbf{VGG-11}: cấu hình A (8 Conv có BN), bỏ ba tầng FC 4.096 nơ-ron thay bằng GAP, \SL{ch3_svhn/vgg11/params} tham số.
  \item \textbf{MobileNetV2}: 17 khối Inverted Residual, \SL{ch3_svhn/mobilenetv2/params} tham số.
  \item \textbf{ResNet-18}: 8 khối Residual, \SL{ch3_svhn/resnet18/params} tham số.
\end{enumerate}
Độ trễ suy luận được đo lại sau khi mọi huấn luyện kết thúc (máy không chạy việc khác): lô 1 ảnh, trung vị của 300 lần
trên GPU và 100 lần trên CPU 4 luồng; thông lượng FPS đo trên GPU với lô 256.

\bang{tab_ch3_acc}{Độ chính xác của các biến thể CNN huấn luyện từ đầu trên SVHN, EuroSAT và 1D CNN trên dữ liệu tiểu đường.}{ch3-acc}
\bang{tab_ch3_cost}{Chi phí tính toán của các biến thể CNN: tham số, FLOPs, độ trễ và thông lượng.}{ch3-cost}

\hinhpdf{f3_curves_svhn}{1}{Mất mát huấn luyện và độ chính xác kiểm thử theo epoch của các kiến trúc trên SVHN.}{curves-svhn}
\hinhpdf{f3_curves_eurosat}{1}{Mất mát huấn luyện và độ chính xác kiểm thử theo epoch của các kiến trúc trên EuroSAT.}{curves-eurosat}
\hinhpdf{f3_cm_svhn}{0.82}{Ma trận nhầm lẫn (\% theo hàng) của ResNet-18 trên tập kiểm thử SVHN.}{cm-svhn}
\hinhpdf{f3_cm_eurosat}{0.95}{Ma trận nhầm lẫn (\% theo hàng) của ResNet-18 trên tập kiểm thử EuroSAT.}{cm-eurosat}
\hinhpdf{f3_perclass_svhn}{1}{F1 theo từng lớp của bốn kiến trúc trên SVHN.}{perclass-svhn}
\hinhpdf{f3_perclass_eurosat}{1}{F1 theo từng lớp của bốn kiến trúc trên EuroSAT.}{perclass-eurosat}
\hinhpdf{f3_wrong_svhn}{1}{Mười sáu ảnh SVHN bị ResNet-18 phân loại sai (nhãn thật → dự đoán, xác suất của lớp dự đoán).}{wrong-svhn}
\hinhpdf{f3_wrong_eurosat}{1}{Mười sáu ảnh EuroSAT bị ResNet-18 phân loại sai.}{wrong-eurosat}
\hinhpdf{f3_latency}{1}{Độ trễ suy luận (GPU, CPU, lô 1), thông lượng GPU (lô 256) và quan hệ độ chính xác -- kích thước
(diện tích điểm tỉ lệ với FLOPs).}{latency}

\paragraph{Phân tích đánh đổi thiết kế.}
\begin{itemize}
  \item \textbf{Độ chính xác.} Trên SVHN, ResNet-18 đạt \SL{ch3_svhn/resnet18/accPct}\,\%, MobileNetV2
    \SL{ch3_svhn/mobilenetv2/accPct}\,\%, Custom CNN \SL{ch3_svhn/custom/accPct}\,\% và VGG-11 \SL{ch3_svhn/vgg11/accPct}\,\%:
    bốn kiến trúc cách nhau chưa tới một điểm phần trăm, vì SVHN $32{\times}32$ không đòi hỏi mạng rất sâu. Trên EuroSAT
    $64{\times}64$, ResNet-18 dẫn đầu rõ hơn với \SL{ch3_eurosat/resnet18/accPct}\,\%, tiếp theo là VGG-11
    \SL{ch3_eurosat/vgg11/accPct}\,\% và MobileNetV2 \SL{ch3_eurosat/mobilenetv2/accPct}\,\% gần như ngang nhau, Custom CNN
    \SL{ch3_eurosat/custom/accPct}\,\%: ảnh lớn hơn, kết cấu phức tạp hơn thì độ sâu và kết nối tắt mới phát huy tác dụng.
    Các đường học (Hình~\ref{fig:curves-svhn}, \ref{fig:curves-eurosat}) dao động mạnh ở giai đoạn đầu khi lịch One-Cycle đẩy
    tốc độ học lên đỉnh, rồi hội tụ khi tốc độ học giảm về cuối.
  \item \textbf{Lỗi còn lại ở đâu.} ResNet-18 sai \SL{extra/ch3/svhn_n_wrong} ảnh SVHN và \SL{extra/ch3/eurosat_n_wrong} ảnh
    EuroSAT. Trên EuroSAT, lỗi tập trung ở các lớp thực vật có kết cấu giống nhau: nhiều nhất là cây lâu năm bị nhầm thành
    thảm cỏ và cây hằng năm, tiếp theo là khu công nghiệp nhầm thành đường cao tốc và sông nhầm thành đường cao tốc (cùng là
    dải dài cắt ngang ảnh) (Hình~\ref{fig:cm-eurosat}, \ref{fig:wrong-eurosat}); các lớp rừng và biển/hồ gần như tuyệt đối
    (Hình~\ref{fig:perclass-eurosat}). Trên SVHN, cặp nhầm nhiều nhất là 3 → 9 và 7 → 1, hai cặp có nét gần nhau; nhiều ảnh sai
    là ảnh mờ, bị che, hoặc có nhiều chữ số mà nhãn chỉ ứng với chữ số giữa (Hình~\ref{fig:wrong-svhn}): một phần không nhỏ là
    lỗi mà người cũng khó phân định.
  \item \textbf{Chi phí tính toán.} MobileNetV2 có \SL{ch3_svhn/mobilenetv2/params} tham số, ít hơn ResNet-18 khoảng 5 lần
    và VGG-11 khoảng 4 lần, FLOPs cũng thấp hơn ResNet-18 nhiều lần (Bảng~\ref{tab:ch3-cost}). Nhưng trên GPU, ở lô 1,
    MobileNetV2 lại \emph{chậm nhất} (\SL{bench/ch3/svhn/mobilenetv2/gpu_ms}\,ms so với \SL{bench/ch3/svhn/resnet18/gpu_ms}\,ms
    của ResNet-18 trên SVHN): nó có hơn 50 tầng nhỏ nối tiếp, mỗi tầng một lần khởi chạy kernel, và tích chập depthwise có
    cường độ tính toán thấp nên GPU không tận dụng được. Trên CPU, lợi thế đảo chiều: MobileNetV2 suy luận một ảnh EuroSAT trong
    \SL{bench/ch3/eurosat/mobilenetv2/cpu_ms}\,ms, ResNet-18 cần \SL{bench/ch3/eurosat/resnet18/cpu_ms}\,ms. Ít FLOPs không đồng
    nghĩa nhanh hơn trên mọi phần cứng; lợi thế của MobileNet thể hiện ở CPU và thiết bị di động, nơi số phép tính mới là nút thắt.
  \item \textbf{VGG-11.} Nhiều tham số thứ hai nhưng kiến trúc phẳng, ít tầng, nên độ trễ GPU thấp
    (\SL{bench/ch3/svhn/vgg11/gpu_ms}\,ms); độ chính xác đứng sau ResNet-18. Đây là minh chứng cho bước tiến của kết nối tắt.
  \item \textbf{Custom CNN.} Chỉ \SL{ch3_svhn/custom/params} tham số, tệp khoảng 1\,MB, nhanh nhất ở mọi phép đo, đủ tốt cho
    SVHN và là mốc kiểm chứng ban đầu hợp lý. Bản Keras cùng kiến trúc, cùng số epoch nhưng huấn luyện bằng Adam (tốc độ học
    $10^{-3}$, giảm theo cosine) và tăng cường dữ liệu bằng các lớp có sẵn của Keras (dịch, lật; không có chuyển vị). Trên SVHN
    nó đạt \SL{ch3_svhn/custom_keras/accPct}\,\%, sát bản PyTorch (\SL{ch3_svhn/custom/accPct}\,\%); trên EuroSAT chỉ đạt
    \SL{ch3_eurosat/custom_keras/accPct}\,\% so với \SL{ch3_eurosat/custom/accPct}\,\%, và đường độ chính xác kiểm thử của nó
    dừng ở mức thấp hơn trong khi các mô hình PyTorch hội tụ dần khi tốc độ học One-Cycle giảm về cuối
    (Hình~\ref{fig:curves-eurosat}). Khoảng cách này đến từ công thức huấn luyện (SGD Nesterov với lịch One-Cycle tổng quát
    hóa tốt hơn Adam trên CNN có BatchNorm, cộng thêm phép chuyển vị), không phải từ nền tảng. Thời gian huấn luyện Keras dài
    hơn nhiều lần vì TensorFlow trên Windows không dùng được GPU.
  \item \textbf{1D CNN trên dữ liệu bảng} đạt F1 \SL{ch3_diabetes1d/cnn1d/f1}, ngang MLP của Chương~\ref{ch:hoc-may}
    (\SL{ch2_mlp_torch/diabetes/f1}): tích chập 1D dùng được, nhưng với dữ liệu bảng thứ tự cột là tùy ý nên “tính cục bộ”
    không mang ý nghĩa như với ảnh, và lợi ích so với MLP là không đáng kể.
\end{itemize}

\subsection{Mạng đã học được gì?}

Hình~\ref{fig:filters} là 64 bộ lọc $3{\times}3$ của tầng thân ResNet-18 huấn luyện trên EuroSAT, mỗi bộ lọc hiển thị như một
ảnh màu nhỏ: nhiều bộ lọc mang một sắc màu chủ đạo (lục, tím, đỏ) tức là nhạy với tỉ lệ giữa các kênh màu, nhiều bộ lọc có
hai nửa tương phản, tức là bộ dò chuyển tiếp sáng -- tối theo một hướng. Đây là những đặc trưng cấp thấp mà mạng tự học
mà không ai lập trình. Hình~\ref{fig:featmaps} đi theo một ảnh khu dân cư qua mạng: ở layer1 ($64{\times}64$) các bản đồ đặc trưng
còn giữ hình dạng đường phố và mái nhà; ở layer2 chúng thô hơn; tới layer4 ($8{\times}8$) chỉ còn vài vùng kích hoạt mạnh,
là biểu diễn trừu tượng mà tầng cuối dùng để phân loại.

\hinhpdf{f3_filters}{1}{64 bộ lọc $3{\times}3$ của tầng thân ResNet-18 (EuroSAT), mỗi bộ lọc vẽ như một ảnh RGB đã chuẩn hóa.}{filters}
\hinhpdf{f3_featmaps}{1}{Bản đồ đặc trưng của một ảnh khu dân cư qua các tầng của ResNet-18: sáu kênh kích hoạt mạnh nhất
ở mỗi tầng.}{featmaps}

% ---------------------------------------------------------------------
\section{Đóng gói triển khai hệ thống Web App thị giác máy tính}

ResNet-18 và MobileNetV2 của cả hai tập dữ liệu được nạp vào phân hệ \emph{Vision AI} của Web App (mục~\ref{sec:web}). Người
dùng chọn một ảnh kiểm thử hoặc tải ảnh lên; hệ thống trả năm lớp có xác suất cao nhất cùng \emph{bản đồ kích hoạt lớp}
(Class Activation Map, CAM) \cite{zhou2016}. Vì mạng kết thúc bằng GAP rồi một tầng tuyến tính, điểm của lớp $c$ là
$\sum_k w^c_k \cdot \mathrm{GAP}(f_k)$; đổi thứ tự lấy tổng và trung bình cho bản đồ $M_c(x, y) = \sum_k w^c_k f_k(x, y)$,
chỉ ra vùng ảnh đã đẩy mô hình về lớp $c$ (Hình~\ref{fig:cam}, Đoạn mã~\ref{code:app__cam_png}).

\hinhtikz{f3_cam}{Cơ chế Class Activation Map: tổ hợp tuyến tính các bản đồ đặc trưng cuối theo trọng số tầng FC của lớp $c$.}{cam}

\maNguon{app__cam_png}{Tính bản đồ kích hoạt lớp CAM từ bản đồ đặc trưng cuối và trọng số tầng FC}

\hinhpdf{f3_cam_eurosat}{1}{CAM của ResNet-18 trên sáu ảnh EuroSAT đúng (hàng trên: ảnh gốc và lớp dự đoán; hàng dưới:
CAM phủ lên ảnh).}{cam-eurosat}
\hinhpdf{f3_cam_svhn}{1}{CAM của ResNet-18 trên sáu ảnh SVHN: vùng nóng nằm đúng trên chữ số được gán nhãn.}{cam-svhn}

\anhUI{ui_vision}{Phân hệ Vision AI: ảnh vệ tinh EuroSAT, bản đồ CAM của lớp dự đoán và xác suất năm lớp hàng đầu của ResNet-18.}{ui-vision}
\anhUI{ui_vision_svhn}{Phân hệ Vision AI với ảnh số nhà SVHN và mô hình MobileNetV2.}{ui-vision-svhn}

Bản đồ CAM tập trung vào đúng chữ số ở giữa ảnh SVHN (bỏ qua chữ số hàng xóm) và vào vùng kết cấu đặc trưng của ảnh vệ tinh
(Hình~\ref{fig:cam-svhn}, \ref{fig:cam-eurosat}), cho thấy mô hình dựa vào đúng bằng chứng thị giác. Suy luận một ảnh kèm CAM
trên CPU của máy chủ mất vài chục mili-giây, đủ cho tương tác thời gian thực.

% =====================================================================
%  CHƯƠNG 4 — MẠNG NƠ-RON HỒI QUY (RNN)
% =====================================================================
\chapter{Mạng Nơ-ron Hồi Quy (RNN)}
\label{ch:rnn}

\section{Cơ sở lý thuyết mạng nơ-ron hồi quy và mô hình hóa chuỗi}

Dữ liệu bảng và ảnh tĩnh ở hai chương trước thỏa giả định các mẫu độc lập cùng phân phối. Dữ liệu chuỗi vi phạm giả định
đó: giá hôm nay phụ thuộc giá hôm qua, việc người dùng bỏ dịch vụ cuối tháng được báo trước bằng hành vi những ngày giữa
tháng. Thứ tự là thông tin. Một tập dữ liệu chuỗi được biểu diễn thành tensor ba chiều $X \in \mathbb{R}^{B\times T\times D}$
với $B$ mẫu, $T$ bước thời gian và $D$ đặc trưng tại mỗi bước (Hình~\ref{fig:tensor4}).

\hinhtikz{f4_tensor}{Tensor chuỗi ba chiều $[B, T, D]$; mỗi lát dọc là vector quan sát $x_t$ tại một bước thời gian.}{tensor4}

\subsection{Kiến trúc Simple RNN và đồ thị tính toán mở rộng theo thời gian}

Mạng hồi quy đơn giản (Elman RNN) giữ một vector trạng thái ẩn $\mathbf{h}_t \in \mathbb{R}^H$ làm bộ nhớ. Tại mỗi bước:
\begin{align}
  \mathbf{h}_t &= \tanh\big(W_{xh}\mathbf{x}_t + W_{hh}\mathbf{h}_{t-1} + \mathbf{b}_h\big), \label{eq:rnn}\\
  \hat y &= W_{hy}\mathbf{h}_T + b_y \quad \text{(dạng many-to-one dùng trong chương này)},
\end{align}
với $W_{xh} \in \mathbb{R}^{H\times D}$, $W_{hh} \in \mathbb{R}^{H\times H}$ được \emph{dùng chung} ở mọi bước. Mở cuộn theo
thời gian (Hình~\ref{fig:unroll}), RNN trở thành một mạng truyền thẳng sâu $T$ tầng có chung trọng số; nhờ chia sẻ trọng số,
số tham số không phụ thuộc độ dài chuỗi.

\hinhtikz{f4_unroll}{Đồ thị tính toán của RNN: dạng cuộn và dạng mở rộng theo thời gian với trọng số dùng chung.}{unroll}

\subsection{Lan truyền ngược qua thời gian (BPTT) và hiện tượng triệt tiêu đạo hàm}

Huấn luyện RNN dùng lan truyền ngược qua thời gian (Backpropagation Through Time): sai số từ mất mát ở bước cuối đi ngược
qua từng bước của đồ thị mở cuộn (Hình~\ref{fig:bptt}). Gradient của mất mát ở bước $T$ theo $W_{hh}$ cộng dồn đóng góp của
mọi bước trước:
\begin{equation}
  \frac{\partial \mathcal{L}_T}{\partial W_{hh}} = \sum_{k=1}^{T}\frac{\partial \mathcal{L}_T}{\partial \mathbf{h}_T}
  \Bigg(\prod_{j=k+1}^{T}\frac{\partial \mathbf{h}_j}{\partial \mathbf{h}_{j-1}}\Bigg)\frac{\partial \mathbf{h}_k}{\partial W_{hh}},
  \qquad
  \frac{\partial \mathbf{h}_j}{\partial \mathbf{h}_{j-1}} = \mathrm{diag}\big(1 - \mathbf{h}_j^2\big)\,W_{hh} .
  \label{eq:bptt}
\end{equation}

\hinhtikz{f4_bptt}{Lan truyền ngược qua thời gian: gradient từ $\mathcal{L}_T$ đi ngược qua chuỗi Jacobian, yếu dần theo
khoảng cách.}{bptt}

Tích $T - k$ ma trận Jacobian quyết định bước $k$ còn ảnh hưởng bao nhiêu \cite{bengio1994}. Vì $|1 - h^2| \le 1$, chuẩn của
tích bị chặn bởi $|\lambda_{\max}(W_{hh})|^{T-k}$:
\begin{itemize}
  \item nếu $|\lambda_{\max}| < 1$: tích giảm theo hàm mũ về 0, gradient của các bước xa \emph{triệt tiêu}, mạng không học
    được phụ thuộc dài hạn;
  \item nếu $|\lambda_{\max}| > 1$: tích tăng theo hàm mũ, gradient \emph{bùng nổ}, trọng số dao động mạnh hoặc tràn số.
\end{itemize}
Hình~\ref{fig:vanish} cho thấy với $|\lambda_{\max}| = 0{,}8$, tín hiệu từ 30 bước trước (độ dài cửa sổ của bài toán chứng
khoán) chỉ còn cỡ $10^{-3}$.

\hinhtikz{f4_vanish}{Bản chất toán học của triệt tiêu và bùng nổ đạo hàm: chuẩn tích chuỗi Jacobian theo số bước lùi.}{vanish}

\subsection{Kỹ thuật cắt tỉa đạo hàm (Gradient Clipping)}

Chống bùng nổ đạo hàm, Pascanu và cộng sự \cite{pascanu2013} đề xuất co gradient khi chuẩn vượt ngưỡng $\theta$:
\begin{equation}
  \mathbf{g} \leftarrow \begin{cases} \mathbf{g}, & \lVert\mathbf{g}\rVert_2 \le \theta,\\[2pt]
    \dfrac{\theta}{\lVert\mathbf{g}\rVert_2}\,\mathbf{g}, & \lVert\mathbf{g}\rVert_2 > \theta. \end{cases}
\end{equation}
Phép co giữ nguyên \emph{hướng} của bước cập nhật, chỉ giới hạn \emph{độ dài} (Hình~\ref{fig:clip}). Mọi mô hình hồi quy
trong chương dùng $\theta = 1$ (\texttt{clip\_grad\_norm\_} trong PyTorch, \texttt{clipnorm} trong Keras). Cắt tỉa không
chữa được triệt tiêu đạo hàm; việc đó cần thay đổi kiến trúc.

\hinhtikz{f4_clip}{Hình học của cắt tỉa đạo hàm: vector $\mathbf{g}$ được co về hình cầu bán kính $\theta$, giữ nguyên hướng.}{clip}

\subsection{Kiến trúc tế bào LSTM và GRU}

\paragraph{LSTM \cite{hochreiter1997}.} LSTM tách bộ nhớ thành trạng thái ô nhớ $\mathbf{c}_t$ và trạng thái ẩn $\mathbf{h}_t$,
điều khiển bằng ba cổng sigmoid nhận $[\mathbf{h}_{t-1}, \mathbf{x}_t]$ (Hình~\ref{fig:lstm}):
\begin{align}
  \text{cổng quên:}\quad & \mathbf{f}_t = \sigma(W_f[\mathbf{h}_{t-1}, \mathbf{x}_t] + \mathbf{b}_f), \\
  \text{cổng vào:}\quad & \mathbf{i}_t = \sigma(W_i[\mathbf{h}_{t-1}, \mathbf{x}_t] + \mathbf{b}_i), \\
  \text{ứng viên:}\quad & \tilde{\mathbf{c}}_t = \tanh(W_c[\mathbf{h}_{t-1}, \mathbf{x}_t] + \mathbf{b}_c), \\
  \text{cập nhật ô nhớ:}\quad & \mathbf{c}_t = \mathbf{f}_t \odot \mathbf{c}_{t-1} + \mathbf{i}_t \odot \tilde{\mathbf{c}}_t, \label{eq:cell}\\
  \text{cổng ra:}\quad & \mathbf{o}_t = \sigma(W_o[\mathbf{h}_{t-1}, \mathbf{x}_t] + \mathbf{b}_o), \\
  \text{trạng thái ẩn:}\quad & \mathbf{h}_t = \mathbf{o}_t \odot \tanh(\mathbf{c}_t).
\end{align}
Phương trình~\eqref{eq:cell} là phép \emph{cộng}: $\partial\mathbf{c}_t/\partial\mathbf{c}_{t-1} = \mathrm{diag}(\mathbf{f}_t)$,
không có ma trận trọng số nhân vào, nên khi cổng quên gần 1 gradient đi qua nhiều bước mà gần như không suy giảm. Cổng quên
quyết định giữ bao nhiêu ký ức cũ, cổng vào quyết định ghi bao nhiêu thông tin mới, cổng ra quyết định để lộ bao nhiêu ô nhớ
ra trạng thái ẩn.

\hinhtikz{f4_lstm}{Kiến trúc tế bào LSTM: đường cộng của ô nhớ (cam) và bốn khối cổng: quên (đỏ), vào (lục), ứng viên (lam),
ra (tím).}{lstm}

\paragraph{GRU \cite{cho2014}.} GRU gộp ô nhớ và trạng thái ẩn, dùng hai cổng (Hình~\ref{fig:gru}):
\begin{align}
  \mathbf{r}_t &= \sigma(W_r[\mathbf{h}_{t-1}, \mathbf{x}_t] + \mathbf{b}_r), \qquad
  \mathbf{z}_t = \sigma(W_z[\mathbf{h}_{t-1}, \mathbf{x}_t] + \mathbf{b}_z), \\
  \mathbf{n}_t &= \tanh(W_n[\mathbf{r}_t \odot \mathbf{h}_{t-1}, \mathbf{x}_t] + \mathbf{b}_n), \qquad
  \mathbf{h}_t = (1 - \mathbf{z}_t) \odot \mathbf{h}_{t-1} + \mathbf{z}_t \odot \mathbf{n}_t .
\end{align}
Trạng thái mới là tổ hợp lồi của trạng thái cũ và ứng viên; với ba khối trọng số thay vì bốn, GRU ít hơn LSTM 25\,\% tham
số ở tầng hồi quy.

\hinhtikz{f4_gru}{Kiến trúc tế bào GRU với cổng đặt lại $r_t$, cổng cập nhật $z_t$ và ứng viên $n_t$.}{gru}

\subsection{Mạng hồi quy hai chiều (BiLSTM) và Recurrent Dropout}

Khi toàn bộ chuỗi đã có sẵn lúc dự đoán (phân loại cả một tháng hành vi), đọc theo cả hai chiều cho nhiều ngữ cảnh hơn.
BiLSTM \cite{schuster1997} chạy một LSTM xuôi $\overrightarrow{\mathbf{h}}_t$ và một LSTM ngược $\overleftarrow{\mathbf{h}}_t$,
rồi ghép $[\overrightarrow{\mathbf{h}}_T; \overleftarrow{\mathbf{h}}_1] \in \mathbb{R}^{2H}$ (Hình~\ref{fig:bilstm}). Lưu ý với
dự báo chuỗi thời gian, chiều ngược vẫn chỉ đọc dữ liệu \emph{trong cửa sổ quá khứ}, không nhìn trước tương lai.

\hinhtikz{f4_bilstm}{Mạng hồi quy hai chiều BiLSTM tổng hợp ngữ cảnh theo cả hai chiều thời gian.}{bilstm}

Dropout thông thường lấy mặt nạ mới ở mỗi bước thời gian, làm nhiễu bộ nhớ của mạng. Gal và Ghahramani \cite{gal2016} chỉ
ra cách đúng là dùng \emph{một} mặt nạ cho mọi bước của cùng một chuỗi (Hình~\ref{fig:dropout}). Trong thực nghiệm, dropout
được đặt giữa hai tầng hồi quy và trước tầng ra, không đặt trong vòng hồi quy, để PyTorch và Keras dùng được nhân cuDNN/oneDNN
tối ưu.

\hinhtikz{f4_dropout}{Dropout thông thường đổi mặt nạ ở mỗi bước; recurrent dropout giữ một mặt nạ cố định cho cả chuỗi.}{dropout}

% ---------------------------------------------------------------------
\section{Khảo sát thực nghiệm và tiền xử lý hai bộ dữ liệu chuỗi}

\subsection{Tập dữ liệu 1: chuỗi giá chứng khoán Việt Nam (VN30, \SL{eda/stock/start} -- \SL{eda/stock/end})}

\begin{itemize}
  \item \textbf{Nguồn:} lịch sử giao dịch trên Sở Giao dịch Chứng khoán TP.~Hồ Chí Minh của bốn mã trụ cột trong rổ VN30,
    tải từ Yahoo Finance bằng thư viện \texttt{yfinance} \cite{yfinance}: FPT (công nghệ -- viễn thông), VCB (ngân hàng),
    HPG (thép) và VIC (bất động sản -- dịch vụ). Giá đã điều chỉnh theo chia tách và cổ tức.
  \item \textbf{Quy mô:} tổng \SL{eda/stock/n_total} phiên: FPT \SL{eda/stock/n/FPT}, VCB \SL{eda/stock/n/VCB},
    HPG \SL{eda/stock/n/HPG}, VIC \SL{eda/stock/n/VIC} phiên. Mỗi phiên có 9 đặc trưng: Open, High, Low, Close, Volume và
    bốn chỉ báo kỹ thuật tự tính -- trung bình trượt 20 và 50 phiên (SMA20, SMA50), chỉ số sức mạnh tương đối 14 phiên theo
    Wilder (RSI-14) và độ biến động 30 phiên (độ lệch chuẩn lợi suất ngày).
  \item \textbf{Bài toán:} từ 30 phiên gần nhất của FPT, dự báo giá đóng cửa phiên kế tiếp.
\end{itemize}

RSI-14 đo tương quan giữa mức tăng và mức giảm gần đây: với $U$, $D$ là trung bình trượt hàm mũ ($\alpha = 1/14$) của mức tăng
và mức giảm giá đóng cửa, $\mathrm{RSI} = 100 - 100/(1 + U/D)$; RSI trên 70 thường coi là “quá mua”, dưới 30 là “quá bán”.

\bang{tab_sample_stock}{Năm phiên giao dịch tiêu biểu của FPT kèm chỉ báo kỹ thuật.}{sample-stock}
\bang{tab_desc_stock}{Thống kê mô tả các đặc trưng của cổ phiếu FPT.}{desc-stock}

\hinhpdf{f4_stock_price}{1}{Giá đóng cửa (nghìn đồng), đường SMA50 và khối lượng giao dịch (cột mờ) của bốn mã VN30.}{stock-price}
\hinhpdf{f4_stock_candle}{1}{Biểu đồ nến 70 phiên cuối của FPT với SMA20, SMA50 và RSI-14 (vùng đỏ: quá mua, vùng xanh: quá bán).}{stock-candle}
\hinhpdf{f4_stock_returns}{1}{Phân phối lợi suất ngày (trên) và độ biến động 30 phiên (dưới) của bốn mã.}{stock-returns}
\hinhpdf{f4_stock_corr}{0.6}{Ma trận tương quan lợi suất ngày giữa bốn mã.}{stock-corr}
\hinhpdf{f4_acf}{1}{ACF của giá đóng cửa và ACF, PACF của lợi suất ngày FPT (vùng tô là khoảng tin cậy 95\,\%).}{acf}

\paragraph{Nhận xét EDA.}
\begin{itemize}
  \item \textbf{Xu hướng và đổi chế độ:} giá FPT tăng khoảng \SL{eda/stock/growth} lần trong giai đoạn khảo sát, đạt đỉnh
    \SL{eda/stock/test_max_close} đồng trong giai đoạn kiểm thử, trong khi mức cao nhất của giai đoạn huấn luyện chỉ
    \SL{eda/stock/train_max_close} đồng (Hình~\ref{fig:stock-price}). VIC tăng vọt từ 2025. Chuỗi giá \emph{không dừng}:
    kiểm định Augmented Dickey--Fuller cho $p = \SL{eda/stock/adf_price_p}$, không bác bỏ được nghiệm đơn vị.
  \item \textbf{Lợi suất gần như không tự tương quan:} lợi suất ngày là chuỗi dừng (ADF $p = \SL{eda/stock/adf_ret_p}$),
    đuôi dày (độ nhọn dư của FPT \SL{eda/stock/ret_kurt/FPT}), và tự tương quan bậc 1 chỉ \SL{eda/stock/acf1_ret}
    (Hình~\ref{fig:acf}). ACF của \emph{giá} giảm rất chậm vì giá hôm nay gần như bằng giá hôm qua, nhưng ACF của
    \emph{lợi suất} gần như nằm trong dải tin cậy: phần dự báo được của biến động ngày mai là rất nhỏ.
  \item \textbf{Biến động theo cụm:} độ biến động 30 phiên tăng vọt thành từng đợt (2018, 2020, 2022, 2025) rồi lắng xuống
    (Hình~\ref{fig:stock-returns}), đặc điểm quen thuộc của chuỗi tài chính.
  \item \textbf{Tương quan giữa các mã} ở mức trung bình ($r$ từ \SL{eda/stock/corr_min} đến \SL{eda/stock/corr_max},
    Hình~\ref{fig:stock-corr}), phản ánh xu hướng chung của thị trường.
\end{itemize}

\hinhtikz{f4_window}{Cơ chế tạo mẫu cửa sổ trượt 30 phiên và chia theo thời gian chống rò rỉ thông tin tương lai.}{window}

\paragraph{Tiền xử lý chống rò rỉ và chống ngoại suy (Hình~\ref{fig:window}).}
\begin{enumerate}
  \item \textbf{Cửa sổ trượt:} cửa sổ $T = 30$ phiên, bước 1, mỗi mẫu $X_i \in \mathbb{R}^{30\times 9}$.
  \item \textbf{Chia theo thời gian, không xáo trộn:} 80\,\% cửa sổ đầu (\SL{eda/stock/n_train} mẫu, đến
    \SL{eda/stock/train_end}) để huấn luyện, trong đó 10\,\% cuối làm kiểm định; 20\,\% cuối (\SL{eda/stock/n_test} mẫu,
    từ \SL{eda/stock/test_start}) để kiểm thử.
  \item \textbf{Chuẩn hóa theo từng cửa sổ:} lần chạy đầu tiên dùng MinMaxScaler khớp trên tập huấn luyện cho $R^2$ âm trên
    tập kiểm thử, vì giá kiểm thử vượt xa mọi mức giá đã thấy và mạng buộc phải ngoại suy. Cách sửa là biểu diễn mọi cột giá
    trong cửa sổ bằng phần trăm chênh so với giá đóng cửa ngày cuối cửa sổ, và dự báo \emph{lợi suất} phiên kế tiếp
    $r_{t+1} = (P_{t+1}/P_t - 1)\cdot 100\,\%$ thay vì mức giá; dự báo giá thu được bằng $\hat P_{t+1} = P_t(1 + \hat r_{t+1})$.
    Khối lượng lấy log rồi chuẩn hóa bằng thống kê của phần huấn luyện, RSI chia 100.
\end{enumerate}

\subsection{Tập dữ liệu 2: quỹ đạo hành vi người dùng KKBox}

\begin{itemize}
  \item \textbf{Nguồn:} cuộc thi WSDM -- KKBox's Churn Prediction Challenge \cite{kkbox2017}: nhật ký nghe nhạc theo ngày của
    người dùng thật của dịch vụ nghe nhạc trực tuyến KKBox (Đài Loan) và nhãn rời bỏ sau khi hết hạn thuê bao.
  \item \textbf{Quy mô:} \SL{eda/churn/n_users} người dùng lấy mẫu phân tầng từ tệp nhật ký 1,4 GB (giữ nguyên tỉ lệ rời bỏ),
    mỗi người là một chuỗi $T = \SL{eda/churn/T}$ ngày của tháng 3/2017 với $D = \SL{eda/churn/D}$ đặc trưng: số bài nghe
    dưới 25\,\%, 25 -- 50\,\%, 50 -- 75\,\%, 75 -- 98,5\,\% và trọn bài; số bài khác nhau; tổng số giây nghe; cờ có hoạt
    động. Đặc trưng đếm và số giây lấy $\log(1 + x)$. Nhãn: 1 nếu người dùng rời bỏ, 0 nếu ở lại.
  \item \textbf{Chia dữ liệu:} 70/15/15 phân tầng: \SL{eda/churn/n_train} huấn luyện, \SL{eda/churn/n_val} kiểm định,
    \SL{eda/churn/n_test} kiểm thử; chuẩn hóa z-score theo từng đặc trưng bằng thống kê tập huấn luyện.
\end{itemize}

\hinhpdf{f4_churn_rate}{0.62}{Tỉ lệ rời bỏ dịch vụ trong mẫu \SL{eda/churn/n_users} người dùng KKBox.}{churn-rate}
\hinhpdf{f4_churn_traj}{1}{Quỹ đạo trung bình theo ngày (vùng tô: khoảng tin cậy 95\,\%) của số giây nghe, số bài khác nhau
và tỉ lệ ngày hoạt động: nhóm rời bỏ đầu tháng nghe nhiều ngang hoặc hơn nhóm ở lại, rồi giảm dần về cuối tháng.}{churn-traj}
\hinhpdf{f4_churn_days}{1}{Phân bố số ngày có hoạt động trong tháng của hai nhóm người dùng.}{churn-days}
\hinhpdf{f4_churn_corr}{0.95}{Ma trận tương quan giữa các đặc trưng hành vi (trung bình 31 ngày) và nhãn rời bỏ.}{churn-corr}

\paragraph{Nhận xét EDA.} Tỉ lệ rời bỏ là \SL{eda/churn/churn_ratePct}\,\%, mất cân bằng mạnh (Hình~\ref{fig:churn-rate}).
Người rời bỏ có trung bình \SL{eda/churn/active_days_churn} ngày hoạt động trong tháng, so với
\SL{eda/churn/active_days_loyal} ngày của người ở lại (Hình~\ref{fig:churn-days}). Điều đáng chú ý là \emph{hình dạng} quỹ đạo
(Hình~\ref{fig:churn-traj}): năm ngày đầu tháng, nhóm rời bỏ có $\log(1 + \text{số giây})$ trung bình
\SL{eda/churn/secs_first5_churn}, còn cao hơn nhóm ở lại (\SL{eda/churn/secs_first5_loyal}); năm ngày cuối tháng nhóm rời bỏ
chỉ còn \SL{eda/churn/secs_last5_churn} trong khi nhóm ở lại giữ \SL{eda/churn/secs_last5_loyal}. Tín hiệu nằm ở xu hướng
theo thời gian, đúng loại thông tin mà mô hình chuỗi có thể khai thác còn đặc trưng tổng hợp cả tháng thì làm mờ đi. Tuy nhiên
tương quan của từng đặc trưng với nhãn đều yếu (tỉ lệ ngày hoạt động $r = \SL{eda/churn/corr_active}$,
Hình~\ref{fig:churn-corr}): hai nhóm chồng lấn nhau rất nhiều, vì quyết định gia hạn còn phụ thuộc giá, khuyến mãi, phương
thức thanh toán, những thông tin không có trong nhật ký nghe. Đây là bài toán khó; mục tiêu thực tế là xếp hạng rủi ro tốt
(ROC-AUC) chứ không phải độ chính xác tuyệt đối.

% ---------------------------------------------------------------------
\section{Hiện thực hóa mã nguồn trên ba nền tảng}

Cả hai nền tảng dùng chung một kiến trúc (Hình~\ref{fig:arch-rnn}): 2 tầng hồi quy 64 đơn vị, dropout 0,2 giữa hai tầng,
lấy trạng thái ẩn cuối của tầng 2 qua Dense 32 -- ReLU -- Dropout -- Dense 1.

\hinhtikz{f4_arch_rnn}{Kiến trúc mạng hồi quy dùng chung cho bốn loại tế bào và hai nền tảng.}{arch-rnn}

\subsection{Tế bào LSTM viết tay bằng NumPy}

Đoạn mã~\ref{code:rnn_scratch__LSTMCellScratch} cài đặt trực tiếp sáu phương trình LSTM với thứ tự khối trọng số
$(i, f, g, o)$ giống PyTorch; lớp \texttt{LSTMRegressorScratch} xếp chồng nhiều tế bào và nạp trọng số từ mô hình PyTorch
đã huấn luyện.

\maNguon{rnn_scratch__LSTMCellScratch}{Tế bào LSTM và lan truyền tiến qua chuỗi viết tay bằng NumPy}

\paragraph{Kiểm chứng.} Nạp trọng số của mô hình LSTM hai tầng đã huấn luyện trên FPT vào bản NumPy và chạy trên toàn bộ
\SL{eda/stock/n_test} cửa sổ kiểm thử: đầu ra lệch PyTorch tối đa \SL{ch4_scratch/lstm_stock/max_abs_diff} (sai số làm tròn
giữa float32 và float64), RMSE dự báo giá \SL{ch4_scratch/lstm_stock/rmse_scratch} nghìn đồng so với
\SL{ch4_scratch/lstm_stock/rmse_torch_cpu} của PyTorch, tốc độ \SL{ch4_scratch/lstm_stock/infer_ms}\,ms mỗi mẫu khi chạy theo
lô. Các phương trình của tế bào LSTM đã được hiện thực đúng.

\subsection{Hiện thực bằng Keras / TensorFlow}

\maNguon{rnn_keras__build_keras_seq_model}{Mô hình hồi quy many-to-one bằng Keras cho cả hồi quy và phân loại}

\subsection{Hiện thực bằng PyTorch}

Một lớp \texttt{SeqModel} phục vụ cả bốn loại tế bào; vòng huấn luyện cắt tỉa chuẩn gradient trước mỗi bước cập nhật.

\maNguon{rnn_torch__SeqModel}{Mô hình hồi quy tổng quát bằng PyTorch: Simple RNN, LSTM, GRU, BiLSTM}
\maNguon{rnn_torch__train_epoch}{Một epoch huấn luyện PyTorch có cắt tỉa chuẩn gradient}

Cấu hình huấn luyện: Adam $10^{-3}$, cắt tỉa $\theta = 1$. Chứng khoán: lô 64, tối đa 100 epoch, dừng sớm sau 15 epoch không
cải thiện, mất mát MSE trên lợi suất. KKBox: lô 512, tối đa 25 epoch, dừng sớm sau 5 epoch, mất mát entropy chéo nhị phân;
ngưỡng phân loại chọn trên tập kiểm định để F1 lớn nhất.

% ---------------------------------------------------------------------
\section{Đánh giá và đối chuẩn các kiến trúc RNN trên hai tập dữ liệu}

\paragraph{Phương pháp.} Bốn kiến trúc được huấn luyện từ đầu trên cùng cách chia dữ liệu. Ngoài các chỉ số chuẩn, bài toán
chứng khoán có thêm mốc tham chiếu \emph{giữ nguyên giá} ($\hat P_{t+1} = P_t$, tương đương dự báo lợi suất bằng 0) và hai
thước đo khó hơn $R^2$ trên giá: $R^2$ trên \emph{lợi suất} và tỉ lệ đoán đúng chiều tăng/giảm.

\subsection{Dự báo chuỗi thời gian cổ phiếu FPT (hồi quy)}

\bang{tab_ch4_stock}{Đối chuẩn dự báo giá đóng cửa FPT trên tập kiểm thử (\SL{eda/stock/test_start} -- \SL{eda/stock/end}).}{ch4-stock}

\hinhpdf{f4_loss}{1}{Đường cong mất mát huấn luyện (nét liền) và kiểm định (nét đứt) của bốn kiến trúc trên hai tập dữ liệu.}{loss4}
\hinhpdf{f4_pred_stock}{1}{Giá thực tế và dự báo trên tập kiểm thử FPT, kèm mốc “giữ nguyên giá”; dưới là 60 phiên cuối.}{pred-stock}
\hinhpdf{f4_ret_scatter}{1}{Lợi suất thực tế và lợi suất LSTM dự báo trên tập kiểm thử (trái) và phân bố của chúng (phải).}{ret-scatter}

\subsection{Phân loại chuỗi hành vi KKBox}

\bang{tab_ch4_churn}{Đối chuẩn phân loại người dùng rời bỏ trên tập kiểm thử KKBox.}{ch4-churn}

\hinhpdf{f4_roc_churn}{1}{Đường ROC và Precision -- Recall của bốn kiến trúc trên tập kiểm thử KKBox; nét chấm là mô hình
ngẫu nhiên.}{roc-churn}
\hinhpdf{f4_prob_churn}{1}{Phân bố xác suất rời bỏ do LSTM dự đoán, tách theo nhãn thật.}{prob-churn}
\haihinh{f4_cm_churn_lstm}{LSTM}{f4_cm_churn_bilstm}{BiLSTM}{Ma trận nhầm lẫn trên tập kiểm thử KKBox tại ngưỡng tối ưu F1.}{cm-churn}

\subsection{Đối chuẩn PyTorch và Keras}

\bang{tab_ch4_fw}{Đối chuẩn PyTorch và Keras với cùng cấu hình mạng.}{ch4-fw}
\hinhpdf{f4_fw}{1}{So sánh PyTorch và Keras: RMSE trên FPT (nét chấm là mốc giữ nguyên giá), F1 trên KKBox, số tham số và độ
trễ CPU.}{fw}

\paragraph{Thảo luận kết quả.}
\begin{itemize}
  \item \textbf{$R^2$ cao trên giá không chứng minh năng lực dự báo.} Cả bốn mạng đạt $R^2 \approx \SL{ch4_torch/stock/lstm/r2}$
    trên giá FPT, nhưng mốc giữ nguyên giá cũng đạt $R^2 = \SL{ch4_torch/stock/naive/r2}$ với RMSE
    \SL{ch4_torch/stock/naive/rmse} nghìn đồng, ngang hoặc tốt hơn mọi mạng (LSTM: \SL{ch4_torch/stock/lstm/rmse}). $R^2$
    trên lợi suất của các mạng xấp xỉ 0 và tỉ lệ đoán đúng chiều chỉ 43 -- 48\,\%, không hơn tung đồng xu. Các đường dự báo
    trùng gần như hoàn toàn với đường thực tế dịch đi một phiên (Hình~\ref{fig:pred-stock}).
  \item \textbf{Mạng học được rằng câu trả lời tốt nhất là “gần 0”.} Hình~\ref{fig:ret-scatter} cho thấy điều này trực tiếp:
    lợi suất LSTM dự báo có độ lệch chuẩn chỉ \SL{extra/ch4/pred_ret_std}\,\%, trong khi lợi suất thực tế dao động với độ lệch
    chuẩn \SL{extra/ch4/true_ret_std}\,\%, và tương quan giữa hai đại lượng là \SL{extra/ch4/ret_corr_lstm}. Điều này khớp với
    EDA (tự tương quan lợi suất gần 0) và với giả thuyết thị trường hiệu quả dạng yếu. Kết luận đúng là: với dữ liệu giá --
    khối lượng thuần túy, không kiến trúc hồi quy nào khai thác được thông tin vượt mốc ngây thơ; muốn có lợi thế cần thêm dữ
    liệu ngoài chuỗi giá (tin tức, báo cáo tài chính, dòng tiền).
  \item \textbf{Khác biệt giữa các tế bào trên chuỗi giá} vì vậy rất nhỏ (Bảng~\ref{tab:ch4-stock}); Simple RNN với cửa sổ 30
    phiên không bị bất lợi rõ rệt vì tín hiệu dài hạn vốn không có. BiLSTM, nhiều tham số nhất, có RMSE kém nhất trên PyTorch:
    thêm năng lực biểu diễn chỉ thêm khả năng khớp nhiễu. Đường mất mát (Hình~\ref{fig:loss4}) cho thấy điều đó: mất mát kiểm
    định trên FPT gần như không giảm, và với BiLSTM còn tăng vọt trong khi mất mát huấn luyện tiếp tục giảm, dấu hiệu điển hình
    của việc ghi nhớ nhiễu; dừng sớm giữ lại trọng số ở những epoch đầu.
  \item \textbf{KKBox: cơ chế cổng có ích nhưng giới hạn nằm ở dữ liệu.} Các tế bào có cổng đạt ROC-AUC và F1 cao hơn Simple RNN
    (Bảng~\ref{tab:ch4-churn}); GRU tốt nhất với ROC-AUC \SL{extra/ch4/GRU/auc} và AP \SL{extra/ch4/GRU/ap}, so với
    \SL{extra/ch4/Simple RNN/auc} và \SL{extra/ch4/Simple RNN/ap} của Simple RNN (Hình~\ref{fig:roc-churn}). Mức tuyệt đối vẫn
    thấp: chỉ từ nhật ký nghe nhạc, người ở lại và người rời bỏ chồng lấn mạnh, đúng như EDA đã cảnh báo. Phân bố xác suất
    (Hình~\ref{fig:prob-churn}) cho thấy xác suất của nhóm rời bỏ chỉ dịch nhẹ sang phải so với nhóm ở lại. Ở ngưỡng tối ưu F1,
    mô hình chỉ bắt được khoảng một phần ba số người rời bỏ (Recall LSTM \SL{ch4_torch/churn/lstm/recall}) mà vẫn kèm nhiều báo
    động giả (Hình~\ref{fig:cm-churn}); nếu dùng để giữ chân khách hàng, danh sách xếp hạng theo xác suất hữu ích hơn một nhãn cứng.
  \item \textbf{PyTorch và Keras} cho độ chính xác cùng mức với cùng cấu hình (Bảng~\ref{tab:ch4-fw}, Hình~\ref{fig:fw});
    khác biệt nhỏ đến từ khởi tạo ngẫu nhiên và thứ tự lô. Số tham số GRU trùng khít (\SL{ch4_torch/stock/gru/params}) vì Keras
    mặc định \texttt{reset\_after=True} cũng dùng hai vector bias như PyTorch; với Simple RNN và LSTM, PyTorch nhiều hơn vài trăm
    tham số do giữ hai bias riêng $b_{ih}$, $b_{hh}$ (tương đương về toán học, chỉ khác cách tham số hóa). PyTorch trên GPU
    huấn luyện nhanh hơn Keras trên CPU nhiều lần; khi suy luận một mẫu trên CPU, độ trễ hai nền tảng cùng bậc mili-giây.
\end{itemize}

% ---------------------------------------------------------------------
\section{Đóng gói triển khai Web App tương tác thời gian thực}

Hai phân hệ cuối của Web App (mục~\ref{sec:web}) dùng các mô hình PyTorch của chương này:
\begin{enumerate}
  \item \textbf{Dự báo giá cổ phiếu FPT:} người dùng chọn kiến trúc và số phiên hiển thị; hệ thống vẽ giá thực tế và dự báo
    trên giai đoạn kiểm thử, dự báo giá phiên kế tiếp sau ngày dữ liệu cuối cùng, và luôn hiển thị RMSE của mốc “giữ nguyên
    giá” cạnh RMSE của mô hình để người dùng không hiểu nhầm năng lực dự báo (Hình~\ref{fig:ui-stock}).
  \item \textbf{Dự đoán người dùng rời bỏ:} lấy một người dùng trong tập kiểm thử (ngẫu nhiên, hoặc chọn nhóm), vẽ quỹ đạo 31
    ngày và trả xác suất rời bỏ của cả bốn kiến trúc so với ngưỡng riêng của từng mô hình (Hình~\ref{fig:ui-churn}).
\end{enumerate}

\maNguon{app__api_churn}{Endpoint dự đoán rời bỏ: chạy bốn mô hình hồi quy trên cùng một chuỗi 31 ngày}

\anhUI{ui_stock}{Phân hệ dự báo giá cổ phiếu FPT: đường giá thực tế và dự báo của LSTM, dự báo phiên kế tiếp, so sánh với
mốc giữ nguyên giá và các chỉ báo hiện tại.}{ui-stock}
\anhUI{ui_stock_gru}{Phân hệ dự báo giá với mô hình GRU và cửa sổ hiển thị toàn bộ giai đoạn kiểm thử.}{ui-stock-gru}
\anhUI{ui_churn}{Phân hệ dự đoán rời bỏ với một người dùng đã rời bỏ: quỹ đạo nghe nhạc 31 ngày và xác suất của bốn kiến trúc.}{ui-churn}
\anhUI{ui_churn_loyal}{Phân hệ dự đoán rời bỏ với một người dùng ở lại.}{ui-churn-loyal}

Mỗi yêu cầu dự báo (bốn mô hình hoặc hàng trăm cửa sổ) được xử lý trong vài chục mili-giây trên CPU.

% =====================================================================
%  KẾT LUẬN
% =====================================================================
\chapter*{Kết Luận}
\addcontentsline{toc}{chapter}{Kết Luận}
\markboth{Kết Luận}{}
\label{ch:ket-luan}

\section*{Kết quả đạt được}

Tiểu luận đã đi trọn một vòng từ lịch sử, lý thuyết, cài đặt, thực nghiệm đến triển khai cho ba họ mô hình trên bảy tập dữ
liệu thật:
\begin{itemize}
  \item \textbf{Lịch sử AI} được hệ thống hóa qua mười sáu mốc, cho thấy mỗi mô hình hiện đại là lời đáp cho một thất bại cụ
    thể trước đó: lan truyền ngược cho giới hạn XOR, học máy thống kê cho tính giòn của hệ chuyên gia, LSTM cho triệt tiêu
    đạo hàm, ResNet cho thoái hóa mạng sâu.
  \item \textbf{Học máy cơ bản và MLP:} trên dữ liệu y tế mất cân bằng, MLP đạt F1 lớp dương \SL{ch2_mlp_torch/diabetes/f1},
    cao nhất trong bảy mô hình; trên giá nhà Việt Nam, Random Forest vẫn nhỉnh nhất ($R^2 = \SL{ch2_sklearn/house/rf/r2}$);
    trên văn bản, MLP chính xác nhất nhưng quá khớp nhanh. MLP cài đặt bằng NumPy thuần đạt kết quả ngang Keras và PyTorch.
  \item \textbf{CNN:} ResNet-18 huấn luyện từ đầu đạt \SL{ch3_svhn/resnet18/accPct}\,\% trên SVHN và
    \SL{ch3_eurosat/resnet18/accPct}\,\% trên EuroSAT; MobileNetV2 đạt độ chính xác sát nút với số tham số ít hơn khoảng
    5 lần, nhưng không nhanh hơn trên GPU ở lô 1. Mạng CNN chạy bằng các hàm NumPy viết tay cho cùng dự đoán với PyTorch.
  \item \textbf{RNN:} trên KKBox, các tế bào có cổng vượt Simple RNN; trên giá FPT, cả bốn kiến trúc chỉ ngang mốc “giữ nguyên
    giá” dù $R^2$ trên giá đạt \SL{ch4_torch/stock/lstm/r2}. LSTM viết tay bằng NumPy tái tạo đúng đầu ra của PyTorch.
  \item \textbf{Triển khai:} một Web App Flask sáu phân hệ phục vụ mọi mô hình trên CPU với thời gian phản hồi vài đến vài
    chục mili-giây, có bản đồ CAM giải thích dự đoán ảnh.
\end{itemize}

\section*{Bài học rút ra}

\begin{enumerate}
  \item \textbf{Chọn mô hình theo dạng dữ liệu, không theo độ “hiện đại”.} Random Forest vẫn là đối thủ khó vượt trên dữ liệu
    bảng; CNN thắng trên ảnh vì cấu trúc không gian; RNN chỉ có lợi khi chuỗi thật sự chứa tín hiệu theo thời gian.
  \item \textbf{Một con số đẹp cần một mốc so sánh.} $R^2 = 0{,}99$ trên giá cổ phiếu hay Accuracy 91,5\,\% trên dữ liệu tiểu
    đường đều đạt được bởi mô hình ngây thơ. Mọi đánh giá trong báo cáo vì vậy luôn đi kèm mốc tham chiếu và chỉ số phù hợp
    với bài toán (F1 lớp hiếm, $R^2$ trên lợi suất).
  \item \textbf{Tiền xử lý quyết định nhiều hơn kiến trúc.} Chuẩn hóa theo mức giá khiến mọi RNN thất bại khi giá vượt vùng
    huấn luyện; chuẩn hóa theo từng cửa sổ sửa được ngay. Ngưỡng phân loại thay đổi F1 nhiều hơn việc đổi mô hình.
  \item \textbf{FLOPs không phải độ trễ.} MobileNetV2 rẻ hơn nhiều lần về phép tính nhưng chậm hơn ResNet-18 trên GPU ở lô 1;
    hiệu năng thực phụ thuộc phần cứng đích.
  \item \textbf{Viết tay để hiểu, dùng thư viện để chạy.} Các bản NumPy giúp kiểm chứng từng công thức; PyTorch trên GPU nhanh
    hơn nhiều bậc khi huấn luyện thật.
\end{enumerate}

\section*{Hướng phát triển}

Các hướng tiếp theo bám sát lộ trình của học phần: thay RNN bằng Transformer cho chuỗi dài; bổ sung dữ liệu ngoài chuỗi giá
(tin tức, báo cáo tài chính) và thông tin giao dịch thuê bao cho bài toán rời bỏ; dùng học chuyển giao với trọng số tiền huấn
luyện cho ảnh; và kết nối các mô hình trong báo cáo với đồ thị tri thức, RAG và hệ đa tác tử ở các chương sau của học phần.


\cleardoublepage
\phantomsection
\addcontentsline{toc}{chapter}{Tài liệu tham khảo}
{\sloppy\printbibliography[title={Tài liệu tham khảo}]}

\end{document}
